\documentclass[11pt]{article}
\usepackage[a4paper,margin=2cm]{geometry}
\usepackage[T1]{fontenc}
\usepackage{lmodern}
\usepackage{amsmath}
\usepackage{xcolor}
\usepackage{enumitem}
\usepackage{array}
\usepackage{tabularx}
\usepackage{longtable}
\usepackage{booktabs}
\usepackage{parskip}
\usepackage[hidelinks]{hyperref}
\hypersetup{pdftitle={Experiment plan: turning the two findings into laws (version 14)}}

\definecolor{soft}{RGB}{110,110,110}
\definecolor{boxbg}{RGB}{244,246,250}
\definecolor{boxline}{RGB}{160,175,200}
\definecolor{warnbg}{RGB}{252,245,235}
\definecolor{warnline}{RGB}{210,160,90}
\definecolor{chgcol}{RGB}{25,115,60}
\definecolor{oldcol}{RGB}{135,135,135}
\newcommand{\gl}[1]{{\color{soft}\itshape\small (#1)}}
\newcommand{\chg}[1]{\,{\color{chgcol}\bfseries\footnotesize[#1]}}
\newcommand{\chgold}[1]{\,{\color{oldcol}\bfseries\footnotesize[#1]}}
\newcommand{\boxit}[1]{\par\noindent\fcolorbox{boxline}{boxbg}{\parbox{\dimexpr\linewidth-2\fboxsep-2\fboxrule}{#1}}\par}
\newcommand{\warnit}[1]{\par\noindent\fcolorbox{warnline}{warnbg}{\parbox{\dimexpr\linewidth-2\fboxsep-2\fboxrule}{#1}}\par}
\setlist{itemsep=2pt,topsep=3pt}
\renewcommand{\arraystretch}{1.25}

\begin{document}

{\LARGE\bfseries Experiment plan: turning the two findings into laws}\\[3pt]
{\large Version 14, frozen on 20 September 2026. It is the first version written after a campaign of the law was run and read. A3's Kafka campaign stopped itself two runs from the end, and its fifteen complete rounds say that load raises the plateau \gl{P3b} and cannot say whether load leaves the cliff where it is \gl{P3a}. The reason is not the number of rounds: at 50\% load this pair's negative rate is 0.15\%, and the position of a cliff in a curve that shallow is not determined by any amount of running. P3a is therefore judged where the plateau is at least 2\%, as P4 already is, and a load below that is reported as out of reach. The rounds rule is also simulated at the plateau and floor the pilot measures, instead of at figures from the machines this work began on. A3's Kafka campaign is reported under version 9, the version it ran under. No run this plan covers starts before it is public (Section~\ref{sec:freeze}).}\\[2pt]
{\color{soft}\small Grey text in brackets is a plain-word explanation. Green tags such as \chg{D14-1} mark what changed since version 13, and grey tags such as \chgold{D13-1}, \chgold{D12-1}, \chgold{D11-1}, \chgold{D10-1}, \chgold{D9-1}, \chgold{D8-1}, \chgold{D7-1}, \chgold{D6-1}, \chgold{D5-2} and \chgold{D3-5} what changed in earlier versions; the ten tables below say what each change is and why. Prices are Azure's list prices as read again on 16 September 2026, and they change.}

\section*{In one paragraph}
\boxit{The paper has three parts, and this plan serves all of them: \textbf{the mechanism} of what is happening, \textbf{that it is happening}, and \textbf{that the industry is not dealing with it}.\chgold{D3-1}\\[3pt]
Both findings rest on old ideas: late clock readings were shown in 2012, and millisecond rounding is in a 1970 HP manual. What could make them laws is prediction. We write down numbers \emph{before} running, on machines and tools we have never touched, and then check them.\\[3pt]
\textbf{Mode 1 (late notes, the mechanism):} we test one specific law, \emph{plateau, cliff, floor}. The machine's scheduler settings predict the trip length at which negative results suddenly stop, and how sharply they stop.\\[3pt]
\textbf{Mode 2 (rounding and throwing away, the industry):} a tool's clock, its rounding and its filter should predict, before it runs, what it will report. Its documentation and issue tracker show whether anyone has dealt with it.\\[3pt]
Version 3 added what the pilot on the first Azure machines showed: a calibration of the added delay in every session, a study of why that delay does not lengthen the trip one-for-one, three tick kernels, and the tools block. Version 4 splits the work into 47 small campaigns, each answering one question in a few hours. It sets the number of repeats from a simulation of each campaign's own test, drops what earlier campaigns already answered, and lets campaigns run at the same time on separate machine pairs without mixing machines inside any test.\chgold{D4-1}\\[3pt]
Version 5 writes down how the campaigns are run. A program judges every run the moment it ends, by rules fixed here, and stops a campaign when the instrument is in doubt. A watch reads every machine pair every five minutes, and every file is fingerprinted from the machine that wrote it to the analysis. No prediction, campaign, number of rounds or decision rule changes.\chgold{D5-2}\\[3pt]
Version 6 corrects what the first runs of version 5 found, and stage 0 starts again. The calibration must now be \emph{known} within 0.3~ms at every step, where version 5 asked every run to lie within 0.1~ms of it, which no machine could meet. The broker's own packet capture checks the added delay, because ping could not decide 0.05~ms. The Redis client acknowledges messages in batches, as the paper's own settings say, and M0 studies bursty traffic, where the departure it explains appears. The other changes protect runs from faults and keep everything measured.\chgold{D6-1}\\[3pt]
Version 8 says which slices are a test of the law on a given pair, and puts one brake above the noise it brakes. A message cannot arrive sooner than the client's own zero-delay trip, and the plateau the law predicts lives at trips below the slice: the first pair's calibration put that floor at 1.964~ms for Kafka, which leaves A1's two smallest slices with no plateau to stand on. Those slices are reported as out of reach, with the trip that excludes them, and the prediction is judged on the slices the pair can reach. The got-it brake now also has to clear the scatter the session's own zero-delay runs showed.\chg{D8-1}\\[3pt]}
\boxit{Version 14 is the first written after a campaign was run and read. A3's Kafka campaign confirms that load raises the plateau, and leaves P3a unresolved -- not for want of rounds but because at 50\% load the negative rate is 0.15\%, and no cliff can be located in a curve that shallow. Neither more rounds nor more messages narrow it, and both were checked against that campaign's own runs.\chg{D14-3} P3a is now judged where the ordinary plateau is at least 2\%, the guard P4 has carried since version 4; simulated at A3's measured levels that lifts it from 9\% power to 98\%, with false confirmations at 1\%.\chg{D14-2} A load below the bar is reported as out of reach, as an unreachable slice already is.\chg{D14-3} The rounds rule is simulated at the plateau and floor the pilot measures rather than at 0.30 and 0.01.\chg{D14-1} The world those simulations draw on now counts negatives out of a run's messages and moves the cliff between runs, because one spread multiplied onto a rate can represent neither.\chg{D14-4}\\[3pt]
Version 13 corrects the instrument and no prediction. The rule judging P8 read the middle of the cliff from $s + 0.4h$; A8 runs $s + 0.5h$, as P8 has said since version 8, so the rule asked for a trip its own campaign never runs and no result could have satisfied it. It confirmed in 0\% of simulated campaigns at every number of rounds, which the rounds rule reported as forty rounds and underpowered. Corrected, P8 asks 6.\chg{D13-1} That rule now separates one that never confirms from one merely weak.\chg{D13-2} P8's ratio is taken only where both clients ran, and says whether the campaign tested P8 at all.\chg{D13-3} Both clients now measure their clock and refuse one coarser than 25~$\mu$s.\chg{D13-4}\\[3pt]
Version 11 finishes what A5's sessions need. The plan gives each core count its own session with its own calibration, because the delay's effect on the trip has to be measured on the machine as the campaign will run it. The baseline is the other half of that same instrument: it is the zero-delay trip every target trip is placed from, and it moves with the core count for the same reasons the calibration does. The plan named the calibration and said nothing about the baseline, so a session at two CPUs would have placed its trips from a floor measured at eight. Each A5 session now measures its own baseline, at the load its campaign runs, before its calibration.\chg{D11-1}\\[3pt]
Version 10 changes no rule. It records that one was misapplied and puts it right. Rounds are set by simulating a campaign at the run-to-run spread its pair's spread pilot measured, and the plan has said since version 4 that the figure is the one measured \emph{on that pair for that backend}. In practice one figure per pair was computed, the median over every setup the pilot ran with both backends mixed together, and that is what set the rounds for A1, A4 and A5's first simulation. Pooling two variances into one is allowed only where the two populations may be assumed to share a variance, and this instrument's own pilots refute that assumption: on the second x86 pair the spread is 0.213 for Kafka and 0.701 for Redis, a factor of 3.3 in the scatter and about 11 in the variance. Their pooled median, 0.457, estimates neither. The harm is not only waste. At the pooled figure A5 would run 8 rounds and its log would claim 86\% power; simulated at the spread Redis actually has, the same campaign reaches about 64\%, below the 80\% the rule demands. Pooling spends too many runs on the quiet backend and certifies the noisy one as powered when it is not.\chg{D10-1}\\[3pt]
Version 9 splits P3 in two and states its first half the way a claim that something did \emph{not} change has to be stated. P3 asked that load raise the plateau \emph{and} leave the cliff where it is, and a prediction is confirmed here only where every part of it is. Simulated at the first pair's own spread, the cliff half held in 97\% of campaigns under the law and the plateau half in 20\%, so P3 held in 18\% --- and the plateau half held in 20\% of campaigns where the law is false as well, because P3's falsifier moves the cliff and not the plateau. Joined together, one half could only subtract from the other. They are now P3a and P3b, each with its own falsifier. P3a is also restated: the whole interval of the cliff's movement must lie inside the band, not merely the number in the middle of it, which a scattered estimate can put inside by luck.\chg{D9-1}\\[3pt]
Version 7 corrects the instrument that judges the delay. Version 6's first shakedown turned a sound pair away on a reading taken by ping, and ping is not the path our messages take: measured side by side, it reads about half a millisecond high, with a tail several times longer, and it shows a difference between the two paths that TCP does not. Every judgement now rests on the broker's own capture, which times the delay where it is added; ping and the paths are recorded and decide nothing. No prediction, and no rule that decides one, changes.\chgold{D7-1}}

\section*{Pocket dictionary}
{\small
\begin{tabularx}{\linewidth}{@{}>{\bfseries}p{3.1cm} X >{\bfseries}p{3.1cm} X@{}}
Run & one measurement, a few minutes of traffic & Round & one pass through every setup of a campaign once \\
Setup & one combination of settings being tested & Session & one machine pair, booted one way, at one core count, on one day \\
Campaign & runs that answer one question in one sitting, with the machines stopped before and after & Machine pair & a driver machine and a broker machine used together \\
Time slice ($s$) & how long the scheduler lets a helper run before switching; 2.8~ms on the first Azure machine & Tick ($h$) & the kernel's heartbeat; 1~ms if HZ=1000, 4~ms if HZ=250, 10~ms if HZ=100 \\
Trip ($D$) & the real time from sending to arrival & Plateau, cliff, floor & the three parts of the predicted curve (Section~\ref{sec:law}) \\
Anchor slice & the 3~ms slice, included in every slice campaign so that drift between days can be measured and removed & Simulated power & the share of made-up campaigns, generated from the law, in which the decision rule confirms the prediction \\
Go-first priority & a setting that makes a helper jump the queue \gl{real-time priority} & Receiver-only delay & extra delay added only on the receiving program's connection \\
Clock offset & making one program's clock read a known amount early or late \gl{libfaketime} & Precision class & low = rounds to milliseconds, medium = microseconds, high = nanoseconds \\
Retention & the share of results a tool keeps & p50, p99 & the median, and the value 99\% of results fall below \\
Pre-register & write the predictions down and freeze them before running & Power & the chance of detecting an effect if it is really there \\
Calibration & measuring how much a known change moves our reading & Equivalence test & a test that two things are the same within a stated margin \gl{two one-sided tests} \\
Manipulation check & confirming that the setting we changed really changed & Bridge session & a short repeat that ties new results to earlier ones \\
Packet capture & a record of network packets with their arrival times & Counterbalanced & arranged so that each kernel comes first, middle and last on different days \\
Lane & one machine pair with its own queues and logs; two lanes never share a machine & Verdict & what a run's checks decide as it ends: it counts, it is repeated, or it stops its campaign \\
Stop rule & a reason for a campaign to stop itself and wait for a person & Shakedown & the instrument checks a new machine pair passes before its first campaign \\
Fingerprint & a SHA-256 code computed from a file; any change to the file changes it & Ledger & the list of a campaign's runs with what happened to each, failures included \\
Void session & a session broken by a fault: kept and reported, never analysed & Batched acknowledgement & the receiving client confirming many messages in one round trip, instead of one round trip per message \\
95\% interval & a range built so that 95 of 100 such ranges would hold the true value & Bursty traffic & messages sent in clusters, as in the football match, instead of evenly \\
Ping \gl{ICMP} & a message type made for testing reachability, which machines may handle differently from ordinary traffic & sockperf & a tool that measures round trips with TCP or UDP, the traffic applications use \\
\end{tabularx}}

\section*{What changed since version 13, and why}
One guard added to a prediction, and the instrument that sets the rounds corrected. A3's Kafka campaign has run and is reported under version 9, the version in force when it ran; the guard below does not reach back to it. A3's Redis campaign, A2, A5, A8 and the tool campaigns have not run.

{\small
\begin{longtable}{@{}p{1.1cm} >{\raggedright\arraybackslash}p{6.1cm} >{\raggedright\arraybackslash}p{6.0cm} >{\raggedright\arraybackslash}p{2.2cm}@{}}
\toprule
\textbf{No.} & \textbf{Change} & \textbf{Reason} & \textbf{Prompted by results?} \\
\midrule
\endhead
D14-1 & The spread pilot \gl{P0} measures the plateau and the floor as well as the spread, per backend and at the anchor slice, and the rounds rule is simulated at them. It has been simulated at a plateau of 0.30 and a floor of 0.01. The pilot reads, at the 3~ms anchor: 0.0346 and 0.0093 for Kafka on the first x86 pair, 0.0568 and 0.0123 on the Arm pair, and 0.0448 and 0.0161 on the second. A queue that can give neither level is refused rather than falling back to the old figures. & Every round count in this plan was set at 0.30 and 0.01, figures from the machines this work began on and from no pair that has run. The measured plateau is five to fifty times lower, and on Kafka the fall is under four times where the simulation assumed thirty. A shallower cliff is harder to locate, so those counts were set for an easier experiment than the one being run. The levels are read per slice as well as per backend: Kafka's floor is 0.0093 at 3~ms and 0.0321 at 1.5~ms, and one median over both is no slice's floor. & \textbf{Yes:} A3's campaign failing to resolve P3a is what sent us to look at what the rounds had been simulated at. The levels themselves are the pilot's, measured before any campaign ran. \\
D14-2 & P3a is judged only where the ordinary plateau is at least 2\%, which is the guard P4 has carried since version 4. Simulated at the levels A3 measured and at its 15 complete rounds: as frozen, P3a confirms in 9\% of campaigns under the law and 0\% where it is false; with the guard, 98\% and 1\%. The plan asks for at least 80\% and at most 5\%. & P3a asks where the cliff is at the lowest load a campaign runs. On this pair that load's plateau is 0.15\%, and a curve falling from 0.0015 to 0.00047 does not fix a cliff position. Resampled from A3's own runs the halfway point there has a 95\% interval 0.625~ms wide, against 0.141 and 0.237 at the other two loads, and it carries 87\% of the variance P3a compares. The guard is not new machinery: it is the condition P4 already puts on the same plateau. The same answer follows without the simulation, from the widths A3 measured: the two remaining loads give a difference known to $\pm$0.167~ms, so under the law the interval runs $-0.167$ to $0.167$ and lies inside the band, while the 0.171~ms the falsifier moves the cliff over those two loads puts it at 0.004 to 0.338 and outside. That matters because the world these simulations are drawn from is quieter than the runs at these two loads \gl{D14-4}, which would flatter the power rather than understate it. & \textbf{Yes:} A3's campaign. It is decided on the simulation above and the identifiability below, not on whether it would rescue that campaign's own answer, and it does not reach back to it. \\
D14-3 & A load whose plateau falls below that bar is reported as out of reach, with the rate that excludes it, and the prediction is judged on the loads the pair can reach --- as version 8 already does for a slice whose plateau lies under the client's own zero-delay trip \gl{D8-1}. The finding behind it is recorded: on the first x86 pair, at 50\% load, the cliff cannot be located, and no amount of running locates it. & Read off A3's own 270 runs without a model. Holding 5, 8, 11 and 15 of its real rounds, the halfway point's interval at 50\% load runs 0.474, 1.111, 1.280 and 0.625~ms --- it does not converge, where 75\% narrows 2.8 times over the same range. Pooling runs so that a point carries the counts a three-times longer run would give moves that interval from 0.624 to 0.626~ms, while the same test at 75\% narrows it 2.4 times on five times the messages, close to the square root counting predicts. The trouble is not precision, which more rounds or longer runs would buy. & \textbf{Yes:} A3's campaign, read back without a model. \\
D14-4 & The made-up world counts. A run draws its negatives from the messages it sent, with an overdispersion, instead of taking a rate multiplied by a lognormal spread; the cliff moves from one run to the next, so the curve's own slope decides how much of that reaches the rate; and the levels may be given per load. & One spread cannot hold what A3 measured. Its scatter at 50\% load is within a fifth of what counting four to seven negatives produces by itself, while at 75 and 88\% the counts are in the hundreds and most of the scatter is the fall's position moving: solving for that wobble gives 0.257~ms at one load and 0.245 at the other, one number for both. The old form also cannot represent the messages a run sends, so it cannot be asked whether a longer run buys what another round buys --- which D14-3 needed answered. Reproducing every point's scatter is not yet within it, and that is recorded rather than claimed: what it does reproduce is the steep part of the fall, where the halfway point is read, and the outcome A3 had. & \textbf{Yes:} A3's campaign is what it was held against. \\
D14-5 & The rounds cannot be set for the tick predictions, and that is stated rather than left to look like a number. The rounds rule simulates one tick at a time, and P2, P2b and P2c compare two kernels. A2's fifteen campaigns run at the floor of 4 rounds until it can. & Found while re-deriving the counts: asked for P2, the rule answers that the prediction never confirms, which is what it now says where a rule cannot be satisfied rather than where it is weak \gl{D13-2}. A2 is the largest block in this plan, 682 runs over 15 campaigns, and its repeats have never come from a simulation. & \textbf{No.} Found in the instrument, with no tick campaign yet run. \\
\bottomrule
\end{longtable}}

\section*{What changed since version 12, and why}
Two corrections to the instrument, and nothing to any prediction. A8 had not run, so everything there was written before the runs it would be judged against. A3 was running under version 9 and was untouched by it; A1, A4 and A7 had finished unanalysed.

{\small
\begin{longtable}{@{}p{1.1cm} >{\raggedright\arraybackslash}p{6.1cm} >{\raggedright\arraybackslash}p{6.0cm} >{\raggedright\arraybackslash}p{2.2cm}@{}}
\toprule
\textbf{No.} & \textbf{Change} & \textbf{Reason} & \textbf{Prompted by results?} \\
\midrule
\endhead
D13-1 & The program that judges P8 read the middle of the cliff from $s + 0.4h$, a trip A8 does not run; it now reads whichever trip across the cliff the campaign carries, nearest the middle first. On the four-trip ladder that is the trip it already read, so no other campaign's judgement moves. A8 runs \textbf{6 rounds}, from P8's corrected simulation at each pair's own Kafka spread \gl{0.250 on the x86 pair, 0.256 on the Arm pair}: confirmed in 88\% and 86\% of campaigns under the law, and in 0\% of campaigns where it is false. That is 288 runs over the two campaigns where all three trips are reachable. & P8 has asked about the rate at $s + 0.5h$ since version 8 \gl{D4-9}, and A8's design has run that trip since the same version. Only the program disagreed, and it disagreed in the direction that cannot fail: a rule that names a trip its own campaign never runs is not a hard test but an unsatisfiable one, and it reported as a weak prediction rather than a broken rule. & \textbf{No.} A simulation of the rule against its own campaign's design, before any A8 run exists. \\
D13-2 & Where a rule confirms in no simulation at all, at any number of rounds, the rounds rule says so instead of answering with the ceiling and the word underpowered. & The two are different things and read alike: 0\% at every step is not a weak test of the law but a rule that cannot be satisfied, and rounds cannot mend it. Reported as underpowered it would put a broken instrument in the plan's table as a weak result, at the ceiling's cost. & \textbf{No.} It follows from D13-1: the defect was visible in the simulation's own output and was read as a number instead. \\
D13-3 & P8's ratio is taken only at the trips both clients ran, and its answer records which those were and whether there were any. Runs a trip only one client reached are kept and reported as before; they are left out of the ratio alone. & The plan takes the ratio ``at matched setups'', and each client is calibrated on its own because the delay's effect on the trip belongs to the client \gl{D4-9}. Their zero-delay floors therefore differ, and a trip can be out of reach for one and not the other \gl{D8-1}. Where the trip lost is the plateau, P8's ratio is undefined, and an undefined ratio was reported as a prediction that did not hold rather than as a campaign that could not test it --- after the campaign had run. & \textbf{No.} Found by placing A8's trips from two calibrations that differ, which is what two clients calibrated separately will give. \\
D13-4 & Both clients measure the clock they are about to stamp with, refuse to start where it steps coarser than 25~$\mu$s, and record what it read in the run's own files. The Java client has refused since it was written, at 100~$\mu$s; the bar is now 25~$\mu$s and the Python client holds it too. & Metrology's rule of ten: an instrument should divide the tolerance being judged into about ten parts, and below four to one it is not a measurement of that tolerance at all. The tolerance is not the cliff's width but the narrowest bound this plan pre-registers --- P2d asks whether the fitted start stays within a quarter of a millisecond, and P3a whether the halfway point's whole interval lies inside 0.25~ms either way. Ten parts of 250~$\mu$s is 25. At 100~$\mu$s the ratio against those two bounds was 2.5 to 1, under the four-to-one floor. Measured: Windows steps in 998{,}600~ns whether the runtime is the JVM or Python, which is the operating system's answer and not the client's; the Linux testbed reads 1{,}000~ns from the JVM and 232~ns from Python. A8 is therefore run on Linux, and the clients enforce that themselves. & \textbf{No.} The clocks of the machines were measured; no campaign was run. \\
\bottomrule
\end{longtable}}

\section*{What changed since version 11, and why}
The tools are read and their predictions fixed. No tool campaign has run, so everything there is written before the runs it will be judged against. On the law's side nothing changed: A3 was running under version 9, and A1, A4 and A7 had finished unanalysed.

{\small
\begin{longtable}{@{}p{1.1cm} >{\raggedright\arraybackslash}p{6.1cm} >{\raggedright\arraybackslash}p{6.0cm} >{\raggedright\arraybackslash}p{2.2cm}@{}}
\toprule
\textbf{No.} & \textbf{Change} & \textbf{Reason} & \textbf{Prompted by results?} \\
\midrule
\endhead
D12-1 & The seven tools that carried a question mark are classified, and the classification is the
prediction T1 to T4 will test. \textbf{Vegeta: high} \gl{monotonic nanoseconds that survive into
the subtraction; nothing dropped, clamped or replaced}. \textbf{valkey-benchmark, memtier\_benchmark,
librdkafka's rdkafka\_performance, YCSB: medium} \gl{microseconds}. \textbf{k6, hey: low}. Each
tool's version is pinned as read. The falsifiable content is stated per tool: that Vegeta's
per-request latencies track our own instrument to within the sub-millisecond and that it never
emits a negative on a live run; that hey's output is quantised in steps of 0.1~ms, so differences
our instrument resolves below that are invisible in it, and that a run of fewer than a hundred
requests prints a 99th percentile of exactly zero; that no memtier percentile reads below
0.008~ms whatever the true latency; and that valkey-benchmark's reported minimum is 0.000~ms for
any reply faster than 8~$\mu$s. & The plan requires a read-first tool to be classified before any
run and states that the classification is itself a prediction \gl{Section~\ref{sec:tools}}. Seven
were unread and are now read. Fixing them here is what makes them predictions rather than notes:
the audit's readings are on the public record before the campaigns that will test them.
& \textbf{No.} Documentation and source, read without running any tool. \\
D12-2 & One earlier classification is recorded as refuted. RabbitMQ PerfTest was classified
``keeps all, negatives included''. It does not keep them: its difference function returns the
size of the difference, so a negative becomes a positive latency of the same magnitude. Under the
cross-machine mode its own documentation prescribes, a consumer clock behind the producer's
reports the clock offset as delay, with nothing in the output to distinguish it. The corrected
reading stands for the runs that follow; the prediction stands as made, and as failed.
& A prediction that fails is reported, not amended away. It also fails in the more damaging
direction than predicted: a visible negative might alert a reader, a sign-flipped one never will.
Its tracker has never carried a report of it -- the three sign and clock terms return nothing
across 215 issues. & \textbf{Yes:} the audit refuted it. The audit reads documentation and source
and no measurement of ours. \\
D12-3 & Two limitations of the audit's own instrument are recorded with its results rather than
repaired. First, the eight search terms are frozen with the plan and two of them earn little:
``clock skew'' returns nothing in projects whose trackers carry extended debates about
unsynchronised processor clocks, because those threads use other words, and ``negative latency''
returns almost nothing anywhere. Had the protocol leaned on those two it would have found nothing
in several projects. The terms were not changed mid-audit and are not changed here. Second, a
tool's handling of odd values can differ between versions -- one tool's guard on its histogram was
removed in January 2025 -- so every reading is recorded against the version read, and a tool
campaign records the version it ran. & An instrument's weaknesses belong in the record beside what
it found. Changing the terms after seeing which ones worked would make the counts unreportable.
& \textbf{Partly:} the weakness of the terms was learned by running them. What they found is
unaffected; what they may have missed cannot be quantified, and that is stated rather than
estimated. \\
\bottomrule
\end{longtable}}

\section*{What changed since version 10, and why}
One thing, and it is about A5's instrument rather than about anything A5 will find. Version 10 was
frozen earlier the same day; A3 is running under version 9 and is untouched by this. A1, A4 and A7
have finished and none has been analysed.

{\small
\begin{longtable}{@{}p{1.1cm} >{\raggedright\arraybackslash}p{6.1cm} >{\raggedright\arraybackslash}p{6.0cm} >{\raggedright\arraybackslash}p{2.2cm}@{}}
\toprule
\textbf{No.} & \textbf{Change} & \textbf{Reason} & \textbf{Prompted by results?} \\
\midrule
\endhead
D11-1 & An A5 session measures its own baseline \gl{B0} at the core count it will run, at the load
its campaign runs, before its own calibration \gl{C0}. Three sessions, six baseline runs each:
eighteen runs in all, about an hour of machine time. The trips of that session's campaign are
placed from that session's own baseline and calibration, and the campaign's design records which
baseline placed it. & The plan already gives each core count its own session with its own
calibration \gl{D4-6}, and the reason it gives is that the delay's effect on the trip must be
measured on the machine as the campaign will run it. A machine with six of its eight CPUs switched
off is not the machine the calibration was taken on --- and it is not the machine the baseline was
taken on either. The baseline is the zero-delay trip that every target trip is placed from, the
floor the law's plateau sits above; it moves with the core count for exactly the reasons the
calibration does. Naming one and not the other left A5's two-CPU session placing its trips from a
floor measured at eight, which puts every trip in the wrong position relative to the cliff the
campaign exists to find. An instrument is calibrated in the state it measures, or it is calibrated
for a different machine. & \textbf{No.} It follows from the plan's own argument for the
calibration, applied to the other half of the same instrument. Nothing measured prompted it and
A5 has not run. \\
\bottomrule
\end{longtable}}

\section*{What changed since version 9, and why}
Nothing in this plan's rules changes here. What changes is a practice that departed from them, and
this section records the departure rather than quietly repairing it. A3 is running under version 9
as this is written; A1, A4 and A7 have finished, and none of the three has been analysed. What
prompted this version is a measurement of the instrument -- each pair's spread pilot, read per
backend instead of pooled -- and a simulation of the decision rule. No result that tests a
prediction was read to reach any of it.

{\small
\begin{longtable}{@{}p{1.1cm} >{\raggedright\arraybackslash}p{6.1cm} >{\raggedright\arraybackslash}p{6.0cm} >{\raggedright\arraybackslash}p{2.2cm}@{}}
\toprule
\textbf{No.} & \textbf{Change} & \textbf{Reason} & \textbf{Prompted by results?} \\
\midrule
\endhead
D10-1 & The spread a campaign is simulated at is the spread its pair's pilot measured for the
backend that campaign runs. One figure per backend, never the two pooled. Where a block's
campaigns run different backends they may now run different numbers of rounds, and each
campaign's log states the spread its own number came from. & The plan has said this since version
4 \gl{D4-2, and the statistics section in every version since}: the simulation uses the spread
measured on that pair \emph{for that backend}. The practice did not. One median was taken over
every setup the pilot ran, both backends together, and used for both. Pooling two variances is
licensed only by assuming the populations share one, and these pilots refute it: 0.250 and 0.423
on the first x86 pair, 0.256 and 0.289 on the Arm pair, 0.213 and 0.701 on the second x86 pair.
The last is a factor of 3.3 in the scatter and about 11 in the variance; their pooled median,
0.457, is an estimate of neither. & \textbf{No.} An instrument measurement from each pair's
spread pilot, and simulations of the decision rule. \\
D10-2 & What the departure cost, campaign by campaign, checked by re-running the rule at each
backend's own spread. \textbf{A1:} its Redis campaigns were re-simulated at 0.602, well above
Redis's true 0.423 on that pair, and 4 rounds still confirmed P1 in 83\% of campaigns under the
law and 0\% where it is false. Its Kafka campaigns ran at a pooled 0.284 where Kafka's own spread
is 0.250, so they are quieter than the number used. \textbf{A4:} its Kafka spread, 0.256, is the
pooled figure exactly, so nothing changes; its Redis campaigns ran at that 0.256 where Redis's own
spread is 0.289, and re-simulated at 0.289 P9 is confirmed in 99\% of campaigns under the law and 0\% where it is false, comfortably above the bar. \textbf{A7:} one
campaign over both backends at one slice, so the pooled figure is the one its analysis pools over
too. & A departure that changes no finished result still has to be shown not to, campaign by
campaign, rather than asserted. Where the check had to assume something it assumed the worse case:
A1's Redis was re-simulated at a spread half again as large as the one measured. & \textbf{No.}
Simulation only; no campaign's data were read. \\
D10-3 & A5's rounds follow from this: 4 for its Kafka campaign, at a spread of 0.213, and 16 for
its Redis campaign, at 0.701. The block's two campaigns therefore run different numbers of rounds,
which D4-2's rule that a block's campaigns share a number does not survive here and is set aside
for A5 \gl{the rule was written for blocks whose campaigns are compared with one another through a
shared anchor; A5's two campaigns are judged separately, one per backend}. & At the pooled figure
both campaigns would run 8: twice what Kafka needs, and half what Redis needs to reach the power
the rule asks for. & \textbf{No.} Simulation only. \\
\bottomrule
\end{longtable}}

\section*{What changed since version 8, and why}
Version 8 was frozen on 18 September 2026 \gl{freeze 04}, and the campaigns it covers began that
night: A1 finished on the first x86 pair, A4 on the Arm pair and A7 on the first pair, all without
a failed run. Nothing they measured prompts anything here. What prompted this version is the
rounds rule itself, run for A3 before A3 ran: it asked for more rounds than its own ceiling
allows and still fell short, and the reason turned out to be the shape of the prediction rather
than the size of the campaign. No measured run was read to reach this change, and none of the
three finished blocks has been analysed.

{\small
\begin{longtable}{@{}p{1.0cm} >{\raggedright\arraybackslash}p{6.2cm} >{\raggedright\arraybackslash}p{6.0cm} >{\raggedright\arraybackslash}p{2.2cm}@{}}
\toprule
\textbf{No.} & \textbf{Change} & \textbf{Reason} & \textbf{Prompted by results?} \\
\midrule
\endhead
D9-1 & P3 becomes two predictions, judged separately. \textbf{P3a:} load leaves the cliff where it
is. \textbf{P3b:} load raises the plateau. Each has its own falsifier: for P3a the halfway point
moving 0.5~ms between 50\% and 88\% load, as before; for P3b a plateau that does not move with
load at all. A3 still runs at the 3~ms slice, three loads and six trips, in two campaigns, one per
backend \gl{D4-4}; its rounds come from P3a. & A prediction here is confirmed only where every part
of it is confirmed, so a prediction of two parts is confirmed about as often as its weaker part
\gl{statisticians call this an intersection--union test, and its power is no more than the least of
its components}. Simulated a thousand times at the first pair's measured spread, the cliff part
held in 97\% of campaigns under the law and the plateau part in 20\%: P3 as written held in 18\%.
And the plateau part held in 20\% of campaigns in the world P3's own falsifier names, because that
falsifier moves the cliff and leaves the plateau alone. A part that answers the same in both worlds
carries no evidence about which world we are in; joined to another by ``and'', it can only
subtract. & \textbf{No.} The rounds rule is a simulation of the decision rule. No measured run was
read. \\
D9-2 & P3a is judged by the whole interval, not by the point estimate: the 95\% interval of the
halfway point's movement between the lowest and the highest load must lie inside $\pm$0.25~ms.
The difference is taken signed, not as a size. & P3a claims that something did \emph{not} change,
and the standard for such a claim is that the whole interval sit inside the band --- two one-sided
tests, as in equivalence testing and in the bioequivalence rules that have used it for decades.
Version 8 asked only that the point estimate be inside the band, which a scattered estimate can
satisfy by luck: simulated under the law, that reading confirmed 97\% of campaigns at four rounds,
where the interval reading confirms 23\%. The looser reading is the more flattering one, which is
reason to distrust it rather than to keep it. Taking the difference signed matters too: folding the
draws about zero pushes the interval's top outwards and would refuse sound campaigns for a reason
that is arithmetic rather than physical. & \textbf{No.} Simulation only. \\
D9-3 & A3 runs 16 rounds, so 576 runs across its two campaigns, in place of the 216
version 8's table estimated. P3b is expected to be reported as underpowered at that number, and
that expectation is written into the campaign's log before it runs. & The rounds are what the
rule returns for P3a at the first pair's own measured spread. P3b would need far more rounds than
the ceiling allows, and the plan's ceiling rule covers exactly that case: the campaign runs, and
the prediction it cannot power is reported as underpowered rather than quietly dropped or quietly
believed. & \textbf{No.} The spread is an instrument measurement from that pair's spread pilot,
taken before any campaign of this block. \\
\bottomrule
\end{longtable}}

\section*{What changed since version 7, and why}
Version 7 was frozen on 18 September 2026 \gl{freeze 03}, and stage 0 ran under it that night on all three pairs: the first pair passed its staircase calibration and its baseline and spread pilots, the Arm pair passed all of stage 0, and the second x86 pair was stopped by a rule halfway through its spread pilot. Nothing the law predicts changes here. What changes is which slices count as a test of it on a given pair --- decided by each pair's own calibration, before any campaign of the block runs --- and where one emergency brake sits. The session stopped on 18 September is reported with its cause, and the pair starts a new one.

{\small
\begin{longtable}{@{}p{1.0cm} >{\raggedright\arraybackslash}p{6.2cm} >{\raggedright\arraybackslash}p{6.0cm} >{\raggedright\arraybackslash}p{2.2cm}@{}}
\toprule
\textbf{No.} & \textbf{Change} & \textbf{Reason} & \textbf{Prompted by results?} \\
\midrule
\endhead
D8-1 & A slice is a test of the cliff on a pair only where that pair's runs can reach the level the halfway point is measured from \gl{$0.9s$}, the level past the cliff \gl{$2(s+h)$} and at least three of the four trips across it. A slice that cannot is reported as out of reach, with the zero-delay trip that excludes it, and is not run. P1 is judged on the slices the pair can reach: the halfway point inside its band in all but one of them, and at least four such slices. P9 asks for all of them, and at least two. The slope rule is unchanged. & A message cannot arrive sooner than the client's own zero-delay trip, and the law's plateau lives at trips below the slice. On 18 September the first pair's calibration measured that floor at 1.964~ms for Kafka and 0.837~ms for Redis. At 0.75~ms Kafka can place 2 of the 8 trips and at 1.5~ms 4 of 8, and neither reaches $0.9s$: there is no plateau to fall from, so those slices cannot test the prediction however the law behaves. Reachability is read from the calibration, an instrument measurement taken before any campaign of the block runs. & \textbf{Yes:} the first pair's calibration of 18 September \\
D8-2 & The brake that stops a session when a run's ``got it'' median moves must clear the instrument's own noise as well as its share of the delay: more than a quarter of the added delay \emph{and} more than the larger of 0.10~ms \gl{the margin P5(c) already gives the got-it median} and three times the scatter of the session's own zero-delay runs. Each session's calibration records that scatter. & On 18 September the second x86 pair's spread pilot stopped on a shift of 0.268~ms where 0.225~ms had been added, a brake set at 0.056~ms. Across four zero-delay runs of that pair's own calibration, with nothing added at all, its Kafka got-it median ran from 1.718 to 2.077~ms \gl{SD 0.183~ms}. A brake below the noise of the thing it brakes stops sound sessions and says nothing about the instrument. & \textbf{Yes:} the stopped session of 18 September and that pair's own calibration \\
D8-3 & The pilot's harness verdict holds the got-it median to the same 0.10~ms and records how far the trip moved per millisecond added, instead of failing when it does not move one-for-one. & The harness replays the match in bursts, and in bursts the trip does not grow one-for-one --- which is why every session's calibration judges that on the campaign's own steady traffic, as version 3 already fixed \gl{D3-5}. The verdict was failing on a number it does not decide, and holding the got-it to half the margin the plan gives it. & \textbf{Yes:} every pair's pilot on 18 September \\
\bottomrule
\end{longtable}}

\section*{What changed from version 6 to version 7, and why}
Version 6 was frozen on 17 September 2026 \gl{freeze 02}, and the first pair started again under it that evening. Its shakedown turned the pair away on the check version 6 had just added, and the measurements that followed showed the check itself was at fault: it read the paths with ping, and ping is not the path our messages take. This version moves every judgement onto the instrument that reads the treatment where it is applied, and leaves ping as a record. Nothing else changes: no prediction, no campaign, no number of rounds, no rule that decides a prediction. Sessions already judged under version 6 keep version 6.

{\small
\begin{longtable}{@{}p{1.0cm} >{\raggedright\arraybackslash}p{6.2cm} >{\raggedright\arraybackslash}p{6.0cm} >{\raggedright\arraybackslash}p{2.2cm}@{}}
\toprule
\textbf{No.} & \textbf{Change} & \textbf{Reason} & \textbf{Prompted by results?} \\
\midrule
\endhead
D7-1 & Every run's added delay --- the number its check is held to, and the number its calibration is fitted at --- is read from the broker's own capture of that run's pings: how much longer it held the replies to the receiver than the replies to the host, within 0.05 ms of the delay set. A run whose capture is missing is repeated. Ping's own figure is written down beside it and decides nothing. & The capture reads the treatment where it is applied: on all three pairs it came within 0.022 ms of the delay set. Ping does not read the path our messages take. Measured side by side on 17 September, over ten blocks in nine minutes, ping read 0.859 and 1.127~ms on the two paths where TCP read 0.447 and 0.611 and UDP 0.439 and 0.695; ping's 90th percentile was 3.2 to 3.7~ms against TCP's 0.46 to 0.62, with single readings of 8.75 and 11.1~ms. Microsoft's own guidance for Azure says ICMP is treated differently from application traffic and names a TCP or UDP tool instead. & \textbf{Yes:} the shakedown of 17 September and the side-by-side probe that followed it \\
D7-2 & With no delay, the two paths are measured by ping, by TCP and by UDP, and recorded. They no longer decide whether a pair may run. What a pair's runs vary by is measured by B0 and sets its rounds, as version 4 already fixed \gl{D4-2}. & On the same probe, ping put the two paths 0.283~ms apart, steadily, where TCP put them 0.009~ms apart: in the receiver's namespace the difference is an artefact of how ICMP and UDP are handled, not something our messages meet. On 17 September that figure turned away a pair whose messages were unaffected --- its trips matched the previous session's within noise --- and a gate on a proxy also chooses which machines contribute data. & \textbf{Yes:} the pair turned away on 17 September, and the probe \\
D7-3 & Every session installs the two measuring tools and records what they read: tcpdump, which times the delay at the broker, and sockperf, whose TCP and UDP round trips on both paths are kept with the machines' facts. & A number that decides nothing still has to be there to explain a pair afterwards. Nothing in the record said what the paths looked like in TCP terms until 17 September, when the question could not be answered without new measurements. & \textbf{Yes:} follows D7-1 and D7-2 \\
D7-4 & The session of 17 September that was turned away is reported as not admitted, with its cause, and the pair was stopped, started again and admitted. Its runs test nothing. & A check that was later found to be the wrong instrument still produced a verdict, and the record has to show both the verdict and the correction. & \textbf{Yes:} that session \\
\bottomrule
\end{longtable}}

\section*{What changed from version 5 to version 6}
Version 5 was frozen on 16 September 2026 \gl{freeze 01}, and stage 0 began that day on three machine pairs. Those runs found faults in the machines, in our code and in two of version 5's own checks, and they were stopped. This version corrects all of them, and stage 0 starts again from the beginning under it. A change prompted by results is exactly what a reader must be able to see, so the table says which result prompted each one. No prediction, and no rule that decides a prediction, changes. The runs of 16 September test nothing; they are kept and reported, and what they showed is the reason for the changes below.

{\small
\begin{longtable}{@{}p{1.0cm} >{\raggedright\arraybackslash}p{6.2cm} >{\raggedright\arraybackslash}p{6.0cm} >{\raggedright\arraybackslash}p{2.2cm}@{}}
\toprule
\textbf{No.} & \textbf{Change} & \textbf{Reason} & \textbf{Prompted by results?} \\
\midrule
\endhead
D6-1 & P5(b)'s second part asks that the calibration be \emph{known} within 0.3~ms at every step: its 95\% interval there, from the scatter of the runs about the line, or within the steps where the step medians are joined, lies within 0.3~ms of it. Version 5 asked every run to lie within 0.1~ms of the calibration. How far each run and each step median lies from it is still reported. & A limit on every run tests how much runs vary, not the calibration, and more runs make it fail more often. On the second x86 pair, run trips scattered around the calibration with a standard deviation of 0.27~ms \gl{Kafka} and 0.18~ms \gl{Redis}. Simulated, 0.1~ms on every run fails every such session, and more than half even at the first pair's 0.06~ms. A check on the step medians alone would pass by itself whenever the calibration \emph{is} the joined step medians, which is when a calibration from few runs is least certain. An interval means the same under both forms. Every run is analysed at the trip it actually had \gl{Section~\ref{sec:stats}}, so this limit sets how closely runs land where they were aimed, not what they show. & \textbf{Yes:} the second x86 pair's calibration failed version 5's check. The first pair's calibration, 51 of its 56 runs, was also fitted before this version was written. \\
D6-2 & C0 runs in two stages on every pair but the first: 2 rounds, then 2 more, fitted together with the first 2, only when the calibration's precision \gl{D6-1} is the one part of its gate that failed. The first pair's staircase S0-1 keeps its 4 rounds. & More rounds only where the data need them, by a rule fixed now. Simulated at the spreads seen: at the first pair's 0.06~ms no session needs more rounds; at 0.18~ms 16\% do, and none then fails; at 0.27~ms 77\% do, and 7\% then fail. A second stage adds about 1.5 hours to a session. & \textbf{Yes:} follows D6-1 and the second pair's spread \\
D6-3 & The shakedown's delay check is made by the broker's own packet capture: at each step, the broker holds the receiver's replies for the set delay beyond the host's, and the host's no longer than with no delay, within 0.05~ms. Ping still checks the delay end to end, within 0.25~ms plus 5\%, the limit every run is held to. & On the Arm pair, one check by ping varied with a standard deviation of 0.03--0.04~ms, and the zero-delay reading drifted by up to 0.1~ms within half a minute. Ping cannot decide 0.05~ms there: the Arm pilot failed the check although the capture shows its delay right to within 0.01~ms. The median of five ping readings would still fail about 8\% of sound Arm shakedowns. The capture times the delay where it is added, to a few microseconds, on both processors \gl{between $-$0.001 and +0.022~ms on all three pairs}. & \textbf{Yes:} the Arm pair's shakedown \\
D6-4 & With no delay, the host's and the receiver's ping round trips must agree within 0.25~ms, as the median of five readings. Each pilot keeps both machines' network and timing facts. A pair that passed its shakedown repeats this network part in every later session. & The receiver-only delay assumes the two paths are the same apart from the delay. Captures on both machines put the gap in Azure's network, not in either machine, and it moved between +0.06 and $-$0.44~ms from one start to the next: a start after a stop can put a machine on another physical host. The limit is the one every run's delay is held to, and it was chosen after the gaps were seen: 0.44~ms on the second x86 pair on 16 September, 0.11 and 0.17~ms on the two x86 pairs on 17 September. & \textbf{Yes:} the second x86 pair's paths, and the probes of 17 September \\
D6-5 & Every run's added delay is measured beyond the same ping taken with no delay just before that run, and the broker's capture of the run's pings is kept beside it. & With no delay, the second x86 pair's paths differed by 0.44~ms, so its zero-delay runs read as $-$0.44~ms of delay and were repeated. A difference that is there without any delay is not delay. The capture gives the delay to the microsecond at the 16 and 32~ms steps, where ping prints only tenths \gl{D6-13}. & \textbf{Yes:} the second pair's first calibration runs \\
D6-6 & Every campaign run uses the paper's own client settings, and each client's log must show that they took, or the run is repeated: Kafka's producer keeps up to 64 requests in flight, as before, and Redis's consumer acknowledges messages in batches of 200, which is new. & The supplement names both settings as learned; the law campaign passed only Kafka's. With one acknowledgement per message, a message waits one round trip for every message read before it. At 16~ms of delay the Redis consumer fell behind to six messages a read and 128~ms trips, which measured its own backlog, not the delay; at 32~ms it would never catch up. On the campaign's steady traffic, up to 8~ms, every read returned one message \gl{6,500 of 6,500}, so nothing changes there. Batches of 200 send one acknowledgement every 4~s, which can make at most one message in 200 wait one extra round trip. This setting failed to apply, silently, three times in earlier work; Kafka's producer now writes its setting down too. & \textbf{Yes:} the first pair's 16~ms staircase step \\
D6-7 & M0 replays the football match in bursts, as the pilot did, and changes the Redis consumer's acknowledgement between one per message and batches of 200. & The departure from one-for-one that M0 exists to explain appears only in bursts. In all five pilots, 2~ms of delay lengthened Redis's trip by 2.5--3.2~ms and Kafka's by 1.7--2.0~ms, while on the campaign's steady traffic both x86 calibrations found slopes of 0.99--1.02 up to 8~ms. In the pilot's replay, 46\% of sends came under 3~ms after the one before, and 11--13\% of Redis reads return several messages, each waiting for the one before it to be acknowledged. M0 as frozen would have run steady traffic and found nothing to explain. Changing the acknowledgement is an intervention on M-H2, the receive loop; M-H1 to M-H4 and their tests are unchanged. & \textbf{Yes:} the calibrations and pilots of 16 September \\
D6-8 & Before every run the driver checks that the receiver's address is only in the receiver's namespace, that the broker answers, and that no package manager is running; a failed check fails the run. Automatic package upgrades are off on every machine, and the receiver's address is removed from the driver's own network settings. & On 16 September an automatic upgrade restarted the first pair's network service during its pilot. The service put the receiver's address back on the driver, the driver lost its broker, and the load read 78.7\% instead of 75\%. The shakedown failed, as it should. & \textbf{Yes:} a fault of the machines in a pilot \\
D6-9 & Each run's folder carries its campaign's name, and a run never writes into a folder that already exists: the campaign stops instead. & Run keys repeat from one campaign to the next. On 16 September a restarted calibration wrote its first run into the folder of the attempt before it, and the watch showed a 35-minute trip. & \textbf{Yes:} a fault of our code \\
D6-10 & A session broken by a fault, or stopped to fix one, before its calibration's gate is \emph{void}: kept, copied, and reported with the cause, never analysed. Everything a driver measured outside a campaign \gl{a pilot, a void session, a probe of the instrument, a log} is copied and fingerprinted like a campaign, and a pair's whole folder of runs is copied before the pair is deleted. Every run also keeps the TCP counters of both machines and the broker's own log. & Four sessions were void on 16 September, two on each x86 pair: one broken by the upgrade of D6-8, one by the zero-delay reading of D6-5, and the two running when the fault of D6-9 was found. Kafka runs on both pairs held pauses of 0.15--0.9~s, in 18 of 54 runs and never more than 39 messages, whose cause no file could show; a TCP segment sent again, or the broker pausing, would show in these files. & \textbf{Yes:} faults of the machines and of our code \\
D6-11 & A failed calibration gate ends the session, and the pair starts a new one, with a new C0. The 20-minute deallocation of an idle pair stays; no session is held open while a rule changes. & Version 5 said only that the main runs do not start. Holding a session open across a rule change would judge its data by a rule written after they were seen. A new C0 and B0 cost about two hours. & \textbf{Yes:} a draft of this version would have carried a running session across the change \\
D6-12 & Stage 0 starts again from the beginning on every pair: the shakedown, then S0-1 on the first x86 pair and C0, B0 and P0 on the others. After any change to the code, the first x86 pair runs alone until its shakedown and first calibration pass, and then the others start. & Every run this plan analyses is judged by rules written before it. On 16 September every fault reached all three pairs at once; a staggered start makes a fault in new code cost one pair a few hours. & \textbf{Yes:} the runs of 16 September \\
D6-13 & Delays measured by ping are known to the precision ping prints: 0.001~ms below 1~ms, 0.01~ms from 1 to 10~ms, and 0.1~ms from 10 to 100~ms. The broker's capture \gl{D6-5} gives the finer value at the 16 and 32~ms steps. & All eight runs of the first pair's 16~ms step read the receiver's round trip as exactly 17.0~ms. & \textbf{Yes:} the first pair's staircase \\
\bottomrule
\end{longtable}}

\section*{What changed from version 4 to version 5}
Version 4 was prepared for freezing but was never frozen, and no run of the campaign has taken place. A change prompted by what a pilot showed is exactly what a reader must be able to see, so every change is listed here and tagged where it appears. Version 5 adds how the campaigns are run, checked, stopped, watched and kept (Section~\ref{sec:operations}). No prediction, campaign, number of rounds or decision rule changes.

{\small
\begin{longtable}{@{}p{1.0cm} >{\raggedright\arraybackslash}p{6.6cm} >{\raggedright\arraybackslash}p{5.9cm} >{\raggedright\arraybackslash}p{1.9cm}@{}}
\toprule
\textbf{No.} & \textbf{Change} & \textbf{Reason} & \textbf{Prompted by pilot results?} \\
\midrule
\endhead
D5-1 & Every new machine pair passes a shakedown before its first campaign: the pilot's scheduler-settings, receiver-only-delay, load and never-negative checks. Its harness-delay check is replaced by the pair's own C0, and its go-first result is recorded but is not a gate (Section~\ref{sec:operations}). & The first pair's pilot caught a false assumption before any run counted \gl{D3-5}. A pair in another region or on another processor gets the same chance to fail before its runs count. Go-first is left to A8, which tests it as a prediction. & Partly: what the first pilot caught \\
D5-2 & Every run is judged as it ends, on its own driver, by fixed limits with three verdicts: it counts; it is repeated, when a condition it was meant to have did not take; or it stops its campaign, when the instrument is in doubt. Two checks are added: at least 99\% of the planned messages sent and arrived, and the added delay within 0.25~ms plus 5\% of its setting. P5(c)'s stop is applied run by run, against the session's zero-delay ``got it'' median. & Version 4 listed the integrity checks but not their limits, when they are applied, or what follows a failure. Limits applied by a program cannot be bent after a result is seen, and they hold when nobody is watching. & No \\
D5-3 & A campaign stops itself when its attempts keep failing: the last 3 all failed, or more than 20\% of the last 20, once 10 have finished. A run is attempted at most 3 times, and a run whose third attempt fails is abandoned and reported. & A machine that has gone wrong at night should cost minutes, not a night of failed runs. & No \\
D5-4 & A finished campaign's runs are copied from its driver, with a SHA-256 fingerprint for every file checked after the copy, before the next campaign starts on that pair. The fingerprints, and a record of the machine, the code version and the time, are kept with the data. A pair is deleted only after all its campaigns are copied. & Results get lost silently in copies. A fingerprint makes a truncated or changed file visible, and copying after every campaign means a lost machine costs at most the campaign in progress. & No \\
D5-5 & A watch outside the testbed reads every pair every 5 minutes, raises alerts, warnings and notices, and deallocates a pair idle for 20 minutes. & Campaigns run unattended, often at night. A stopped campaign or a broken machine is seen within minutes, and idle machines stop billing. & No \\
D5-6 & Every run records its pair, its driver and the version of the code \gl{the commit} that ran it. All pairs run the same version, and a campaign's queue runs on one pair only. & D4-10 records the pair. Recording the code version too ties every result to the exact code behind it, and one version on every pair keeps the pairs comparable. & No \\
D5-7 & All three machine pairs run in Sweden Central, so the second x86 pair does not go to Italy North. The subscription's limits there were raised on 16 September 2026 to 40 CPUs in total, 20 in the x86 family and 10 in the Arm family. & One region is cheaper \gl{EUR 0.4165 an hour a pair against 0.4577}, puts both x86 pairs on the same hardware generation, and lets all three pairs run at once, about 5 days. The prices were read again on 16 September and had not moved. & No \\
\bottomrule
\end{longtable}}

\section*{What changed from version 3 to version 4}
Version 3 was approved but not frozen.

{\small
\begin{longtable}{@{}p{1.0cm} >{\raggedright\arraybackslash}p{6.6cm} >{\raggedright\arraybackslash}p{5.9cm} >{\raggedright\arraybackslash}p{1.9cm}@{}}
\toprule
\textbf{No.} & \textbf{Change} & \textbf{Reason} & \textbf{Prompted by pilot results?} \\
\midrule
\endhead
D4-1 & The work is split into 47 small campaigns, each answering one question in one sitting \gl{3.5 to 8.4 hours at the floor of 4 rounds}, with the machines stopped between them (Section~\ref{sec:campaigns}). Version 3's two stages are replaced. & A design decision, 15 September: results arrive in pieces that can each be checked, and no machine waits idle between blocks. & No \\
D4-2 & Rounds are set for each campaign by simulating its own prediction and decision rule at the run-to-run spread its machine pair measured, with a floor of 4 and a ceiling of 40. Tool campaigns run 4 rounds. This replaces version 3's fixed floors of 15 and 10 (Section~\ref{sec:stats}). & The fixed floors came from old setups whose spread was about as large as the rate itself. The pilot's repeat runs varied far less, and simulating the actual decision rule measures power directly. & Yes: the pilot's run-to-run spread \\
D4-3 & A1 runs as six campaigns, three per backend, each including the 3~ms slice as a shared anchor. Each campaign's offset is removed through the anchor before slices are compared. & The slices no longer share one session; the anchor keeps the comparison fair across days. & No \\
D4-4 & A3 tests one slice \gl{3~ms} instead of two, in two campaigns, one per backend. P3 is judged at that slice. & Earlier campaigns \gl{E-A3, E-A6} already showed load raising the negative rate. What is new is whether load moves the cliff, and one slice tests that. & No \\
D4-5 & A7 tests one slice \gl{3~ms}, both backends, in one campaign. P4 is judged at that slice. & Go-first priority was shown on the earlier machines \gl{E-A5} and again, 27-fold, in the pilot on this machine. What is new is its effect at the cliff. & Partly: the pilot reproduced go-first \\
D4-6 & A5 runs as two campaigns, one per backend. Each core count is its own session with its own calibration. & Switching CPUs off changes the machine, so the delay's effect is measured again. & No \\
D4-7 & A2 runs as a build campaign, twelve sessions \gl{kernel $\times$ day $\times$ backend} and two bridge sessions on the stock kernel. & Sessions of about 3.5 hours, with each kernel on two days in a balanced order. & No \\
D4-8 & A4 opens with its own instrument campaign on the Arm pair \gl{C0, B0 and P0}, then four campaigns: two per backend, at slices of 1.5 and 3~ms, and of 3 and 4.5~ms. & The repeats rule needs each pair's own spread \gl{D4-2, D4-10}. & No \\
D4-9 & A8 is kept, narrowed to Kafka, whose official client is written in Java, in two campaigns: one on an x86 pair and one on the Arm pair. P8 says how the cliff is judged at three trips. & The earlier client comparison \gl{M1} tested a second Python client on a different delay, so it does not answer P8. & No \\
D4-10 & Several machine pairs may run at the same time. Every prediction is tested on runs from one pair, every pair opens with its own B0 and P0, and every run records its pair. & A shorter calendar, without mixing machines inside any test. & No \\
D4-11 & Straight reruns of earlier campaigns are dropped. A new section lists what earlier campaigns already answered and what is still run (Section~\ref{sec:earlier}). & Nothing is repeated that a new prediction does not need. & Partly: the pilot reproduced go-first \\
D4-12 & Stage 0 opens with a staircase campaign \gl{S0-1}: added delays up to 16~ms, both backends, 4 rounds, before any other campaign. & A two-round session calibration has little power to test whether the line is straight; this campaign has more, and it shows what later calibrations should look like. & Yes: the pilot's delay result \\
D4-13 & The plan is frozen by committing it in full to the public repository before the first Stage 0 run; a later change is a new, numbered freeze (Section~\ref{sec:freeze}). & Version 3 listed three ways and left the choice open. & No \\
D4-14 & P9 is restated for what A4 measures: three slices at one load, with the anchor rule of A1. & Version 3's P9 also named P2 and P3, which A4 never tested. A prediction must match its experiment. & No \\
\bottomrule
\end{longtable}}

\section*{What changed from version 2 to version 3}
{\small
\begin{longtable}{@{}p{1.0cm} >{\raggedright\arraybackslash}p{6.6cm} >{\raggedright\arraybackslash}p{5.9cm} >{\raggedright\arraybackslash}p{1.9cm}@{}}
\toprule
\textbf{No.} & \textbf{Change} & \textbf{Reason} & \textbf{Prompted by pilot results?} \\
\midrule
\endhead
D3-1 & The paper's three parts frame the plan, and the tools block \gl{T1--T4} is in this paper. & A design decision, 15 September. & No \\
D3-2 & Azure only. A4 compares Azure's Arm machine with its x86 machine; the Oracle options are dropped. & A design decision, 14 September. & No \\
D3-3 & Default slices are predicted from the per-CPU constant read off each machine, not from the kernel's version \gl{A5, P7}. & The Azure 6.8 kernel carries the constant Linux 6.15 adopted: 2.8~ms at 8 CPUs, where the 6.8 rule gives 3~ms. A reading of the machine's settings, not an experimental result. & No \\
D3-4 & New section: what the pilot showed. & A record of the pilot. & -- \\
D3-5 & Version 2's built-in delay check is replaced by a calibration in every session \gl{C0}. The ``got it'' check becomes an equivalence test. The law is analysed at each run's measured trip. P5 is rewritten. & The pilot showed the trip does not grow one-for-one with the added delay, and in opposite directions for the two clients. & \textbf{Yes} \\
D3-6 & New study M0: why the trip does not grow one-for-one, with four explanations and the test that decides each. & The same pilot result. A law's measuring setup has to be understood, not only corrected. & \textbf{Yes} \\
D3-7 & A2 uses three kernels \gl{ticks of 1, 4 and 10~ms} built by us from one source, with a bridge to the stock kernel, checks at every boot and counterbalanced sessions. P2b, P2c and P2d are added. & Three levels test a proportion, not only a direction. & No \\
D3-8 & New T5: an audit of each tool's documentation and issue tracker, and prediction T-P8. & It turns ``the industry is not dealing with it'' into a measurement. & No \\
D3-9 & Size and schedule corrected: each backend counts as its own runs, and a run takes 3.5~minutes as the pilot measured, not 1.5. & Version 2 underestimated the machine time about four-fold. & Pilot timing only \\
D3-10 & How the plan is frozen is specified. & Version 2 said ``committed with a date and hash'' without saying where. Settled by D4-13. & No \\
\bottomrule
\end{longtable}}

\section[The three parts, and the questions they ask]{The three parts, and the questions they ask\protect\chgold{D3-1}}
\begin{enumerate}
  \item \textbf{The mechanism (Mode 1).} Is there a law that predicts, from the machine's scheduler settings alone, for which trip lengths a latency measurement goes negative? Section~\ref{sec:law} states the one we test.
  \item \textbf{That it is happening.} Does the effect appear on machines, processors and kernels it was never seen on? Every run of this campaign measures it on new hardware, and A4 moves to a different processor family.
  \item \textbf{That the industry is not dealing with it (Mode 2).} Does a tool's source code predict what it reports? We check whether it rounds, whether it averages over negatives, whether it drops or replaces odd values, and whether it measures the whole trip at all. We test this on every free, widely used tool we can run, not just the one we audited. And we check whether its documentation says so, and what happened when users reported it.
\end{enumerate}

\section{The law we are testing for Mode 1: plateau, cliff, floor}\label{sec:law}

\textbf{The idea in plain words.} Most of the time the helper that writes down ``got it'' gets a core almost at once. But sometimes another task is holding the core. The scheduler lets that task finish its protected time slice, plus up to one tick, before the helper runs. The kernel's own developers state this bound: the wait stays within ``slice plus tick''. So a kept-waiting helper waits roughly between one slice and one slice plus a tick. That gives three parts:
\begin{itemize}
  \item \textbf{Plateau.} If the trip is shorter than the slice, almost every long wait outlasts it, and the negative rate stays roughly flat.
  \item \textbf{Cliff.} Between one slice and one slice plus a tick, fewer and fewer waits outlast the trip, so the rate falls roughly in a straight line.
  \item \textbf{Floor.} For trips longer than one slice plus a tick, long waits almost never outlast the trip. Only a small floor remains, from ordinary short delays.
\end{itemize}

\textbf{The same thing as a formula.} Write $s$ for the time slice, $h$ for the tick, $p$ for how often the helper is kept waiting, and $c$ for the floor. Then
\[
\text{negative rate}(D) \;\approx\; c \;+\; p \times
\begin{cases}
1 & D \le s \quad\text{(plateau)}\\[2pt]
\dfrac{s + h - D}{h} & s < D < s + h \quad\text{(cliff)}\\[6pt]
0 & D \ge s + h \quad\text{(floor)}
\end{cases}
\]

\textbf{What the settings predict, and what we still have to measure.}
\begin{itemize}
  \item \textbf{Predicted from settings alone:} \emph{where} the cliff starts, at $s$, and \emph{how wide} it is, $h$. Both can be read from the machine, or set by hand, before any run.
  \item \textbf{Measured, not predicted:} \emph{how high} the plateau is, $p$, which rises with load, and the floor $c$. Very short trips, under about 0.5~ms, also pick up short waits the formula ignores. So the tests below focus on where the cliff sits and how wide it is.
\end{itemize}

\textbf{Does our old data fit?} On the old 8-core machines ($s = 3$~ms, $h = 1$~ms), 13\% of traced waits were over 2~ms but only 2.8\% over 4~ms. That is the cliff, between 3 and 4~ms. It fits, but it is one machine, and a law has to survive new ones. What the first Azure pilot saw is in Section~\ref{sec:pilot}.

\textbf{How it could fail.} It fails if the cliff does not move when we change the slice (P1). It fails if its width does not follow the tick (P2 to P2c). And it fails if load moves the cliff rather than just raising the plateau (P3).

\section[Machines and money]{Machines and money\protect\chgold{D3-2}\protect\chgold{D4-10}}
{\small
\begin{tabularx}{\linewidth}{@{}p{2.8cm} X p{3.1cm}@{}}
\toprule
\textbf{Pair} & \textbf{Machines and region} & \textbf{Price while running} \\
\midrule
First x86 pair & Driver Standard\_D8as\_v6 \gl{8 CPUs, AMD EPYC Genoa} and broker D2as\_v6 \gl{2 CPUs}, in Sweden Central. The pilot ran here. & about EUR 0.42 an hour \\
Second x86 pair & The same sizes, in Sweden Central beside the first pair, in its own resource group and network.\chgold{D5-7} & about EUR 0.42 an hour \\
Arm pair & Driver D8ps\_v6 \gl{8 CPUs, Azure Cobalt 100} and broker D2ps\_v6, in Sweden Central, the only region that offers these sizes to this account. For A4 and A8. & about EUR 0.32 an hour \\
\bottomrule
\end{tabularx}}

All pairs live in one pay-as-you-go Azure subscription. Each region allows a set number of CPUs, and a stopped machine still counts against that limit. On 16 September 2026, on request, Sweden Central's limits were raised to 40 CPUs in total, 20 in the x86 family and 10 in the Arm family, so all three pairs fit in that one region.\chgold{D5-7} A stopped machine costs only its disk and its fixed address, about EUR 26 a month per pair. A pair that will not be used for more than about two weeks is deleted and later recreated, which takes about 15 minutes.

\textbf{Every tool below is open source and free.} Each tool's licence is recorded when it is downloaded.

\section{The tools}\label{sec:tools}
Precision classes come from reading each tool's code. Tools marked ``read first'' are classified before any run, and that classification is itself a prediction (Section~\ref{sec:predictions}). Each tool's version is pinned by release or commit before any run.

{\small
\begin{longtable}{@{}>{\raggedright\arraybackslash}p{3.6cm} >{\raggedright\arraybackslash}p{2.6cm} >{\raggedright\arraybackslash}p{1.5cm} >{\raggedright\arraybackslash}p{5.0cm} >{\raggedright\arraybackslash}p{2.4cm}@{}}
\toprule
\textbf{Tool} & \textbf{Times} & \textbf{Class} & \textbf{Odd values (zero or negative), from source} & \textbf{Runs on} \\
\midrule
\endhead
OpenMessaging Benchmark & send $\to$ receive; send $\to$ ``got it'' & Low; High & end-to-end: drops $\le 0$, uncounted. Publish: keeps all & Kafka, Valkey, RabbitMQ, NATS \\
Kafka EndToEndLatency & send $\to$ receive & Low (chops) & keeps all; average exact & Kafka \\
Kafka ProducerPerformance & send $\to$ ``got it'' & Low & keeps all & Kafka \\
Apache Pulsar perf & send $\to$ receive & Low & drops negatives, keeps zeros, uncounted & Pulsar \\
emqtt-bench & send $\to$ receive & Low & keeps all & EMQX or Mosquitto \\
wrk2 & request $\to$ response & Medium & negatives never recorded; may crash & nginx \\
RabbitMQ PerfTest & send $\to$ receive & High & keeps all, negatives included & RabbitMQ \\
NATS CLI (latency) & send $\to$ receive & High & keeps all & NATS server \\
\midrule
\multicolumn{5}{@{}l}{\emph{Read first, then run (widely used, class not yet known)}} \\
nats bench & to confirm & ? & to confirm & NATS server \\
valkey-benchmark \gl{free Redis fork} & request $\to$ reply & ? & to confirm & Valkey \\
memtier\_benchmark & request $\to$ reply & ? & to confirm & Valkey \\
librdkafka rdkafka\_performance & send $\to$ ``got it''? & ? & to confirm & Kafka \\
YCSB, k6, Vegeta, hey & request $\to$ reply & ? & to confirm & nginx or Valkey \\
\midrule
\multicolumn{5}{@{}l}{\emph{Kept as source readings only (they do not time a message trip)}} \\
fio, btt, Rezolus & disk and scheduler timings & High & fio replaces negatives with 0; btt with 1~ns; Rezolus counts its discards & -- \\
\bottomrule
\end{longtable}}

\textbf{Our own program} is the reference instrument everywhere. It reads nanoseconds, keeps everything, and records all three times.

\section[What the pilot showed]{What the pilot showed\protect\chgold{D3-4}}\label{sec:pilot}
The first pilot ran on 14 and 15 September on the first x86 pair, at 75\% load. A pilot tests the instruments, never a prediction.

{\small
\begin{tabularx}{\linewidth}{@{}>{\raggedright\arraybackslash}p{3.0cm} >{\raggedright\arraybackslash}X >{\raggedright\arraybackslash}p{6.2cm}@{}}
\toprule
\textbf{Check} & \textbf{Rule written before the pilot} & \textbf{Result} \\
\midrule
Scheduler settings & a slice set by hand is still set 60~s later, and the default comes back & Passed: 1.5~ms held; the 2.8~ms default restored \\
Receiver-only delay, by ping & the set delay reaches the receiver's address only & Passed: 0.5~ms measured as 0.515~ms, 2.0~ms as 2.03~ms \\
Load & within 3 points of 75\% & Passed: 75.3\% in all three cells \\
Never negative & arrival never before sending & Passed: none in 18 runs, about 20,000 messages \\
Delay seen by the harness & the trip grows by the added 2.03~ms within 0.05~ms, and ``got it'' moves under 0.05~ms & \textbf{Failed.} Kafka's trip grew 1.80~ms and ``got it'' moved $-$0.14~ms. Redis's trip grew 2.51~ms and ``got it'' moved $-$0.02~ms. \\
Go-first priority & cuts the negative rate at least 5-fold, intervals apart, at the same load & Passed: 13.5\% to 0.50\% \gl{27-fold}, at 75.3\% and 75.4\% load \\
\bottomrule
\end{tabularx}}

\textbf{What the failed check means.} The delay does reach only the receiver: ping says so, and ``got it'' barely moves. But the trip does not grow one-for-one with it. Repeat runs agree within about 0.1~ms \gl{Kafka's trips were 5.29--5.40~ms with the delay and 3.48--3.92~ms without; Redis's 4.56--4.72 and 2.06--2.33}, so this is a property of the two clients, in opposite directions, and not noise. Version 2's rule was an assumption, and it is false. Trips placed by ``delay = target minus baseline'' would miss by up to half a millisecond, as much as the detail of the cliff we want to see. Hence C0 and M0 in Section~\ref{sec:mode1}.\chgold{D3-5}\chgold{D3-6} The pilot replayed the football match at ten times its speed, in bursts. On the campaigns' steady 50 messages a second, the calibrations of 16 September found both clients tracking the delay within 2\% up to 8~ms, so the departure belongs to bursty traffic, and M0 studies it there.\chgold{D6-7}

\textbf{The machine's own settings.} Kernel 6.8.0-1064-azure, tick 1~ms \gl{HZ=1000}. The default slice was 2.8~ms on the 8-CPU driver and 1.4~ms on the 2-CPU broker. That is the per-CPU constant of Linux 6.15, not the 3 and 1.5~ms the 6.8 rule gives.\chgold{D3-3}

\textbf{Consistent with the law, but not a test.} Without the delay, Kafka's trip \gl{about 3.5~ms} sat inside the cliff the law predicts for this machine \gl{2.8 to 3.8~ms}, and Redis's \gl{about 2.2~ms} sat on the plateau. With 2~ms added, both trips passed the start of the floor, and negative results fell from 7--12\% of messages to 1--1.5\% for both clients. That is the direction the law predicts. But pilot data never tests a prediction, and two points are not a curve.

\section[What earlier campaigns already answered]{What earlier campaigns already answered\protect\chgold{D4-11}}\label{sec:earlier}
Our earlier campaigns ran on Oracle Cloud machines of the same AMD generation. Where one of them already answers a question, it is not run again, and only what a new prediction needs is run.

{\small
\begin{longtable}{@{}>{\raggedright\arraybackslash}p{2.9cm} >{\raggedright\arraybackslash}p{6.5cm} >{\raggedright\arraybackslash}p{6.4cm}@{}}
\toprule
\textbf{Earlier campaign} & \textbf{What it showed} & \textbf{What this plan still runs} \\
\midrule
\endhead
E-A5, and the Azure pilot & Go-first priority on the timestamping helper cuts negative results many times over at equal load; 27-fold in the pilot. & Only its effect at the cliff: A7, at one slice. \\
E-A3 and E-A6 & Negative results rise with load; E-A6 also tested whether it matters where on the machine the load sits, at equal utilisation. & Only whether load moves the cliff: A3, at one slice. \\
E-A9 & Helper waits traced from the kernel's scheduler events predicted the negative rate, with ratios of prediction to measurement between 0.78 and 1.32 in ordinary runs. & Rides along in half of the A1 and A3 runs at no extra cost \gl{A6, P6}, on new hardware. \\
E-A10 & Lengthening the true trip by padding messages lowered the negative rate \gl{4.1-fold for a 77-fold longer trip} without touching the scheduler. & The slice campaigns \gl{A1} place the cliff; the padding sweep is not rerun. \\
E-X to E-X4 & Four audits of the OpenMessaging Benchmark: its end-to-end result drops values at or below zero without counting them. & No new runs of that tool; its source reading stays in the tools table. \\
M1 & A second Kafka client for Python \gl{confluent-kafka}, compared on a different delay: the producer's scheduling lag at real-time replay. & It does not answer P8, so A8 compares Python with Java at the cliff. \\
Earlier mechanism campaigns, rerun unchanged & -- & Not run. The pilot already reproduced go-first on this machine, and every campaign here measures the effect on new hardware. \\
\bottomrule
\end{longtable}}

\section{Experiments for Mode 1 (the mechanism)}\label{sec:mode1}

\warnit{\textbf{The fix: slow only the receiver.} Our sender and receiver sit on the same machine, so the broker's replies to both travel the same link. A delay added to that whole link slows the ``got it'' reply as much as the message. Both times then move together and the test sees nothing. That is exactly why our old delay sweep failed. So:
\begin{enumerate}
  \item Give the receiving program its own network address on the client machine \gl{a second IP address, or its own network namespace}, and make it connect from that address only.
  \item On the broker machine, add the delay only to traffic going to that address. Replies to the sender's address are untouched.
  \item \textbf{Checked, not assumed:} the delay's effect on the trip is measured in every session (C0 below), and the ``got it'' delay must stay equivalent (P5). The added delay itself is measured on its own at every step, by the broker's own capture of those pings, where the delay is added, and the measured value is used, not the setting.\chgold{D3-5}\chgold{D6-5}\chgold{D7-1}
\end{enumerate}}

\textbf{Always the same:} background load is measured, not just set. Before every run we read back the machine's actual time slice \gl{the kernel's base-slice setting} and tick \gl{CONFIG\_HZ}, and record its core count, CPU model, kernel version and machine pair.\chgold{D4-10}

\textbf{C0. Calibrating the added delay, in every session.}\chgold{D3-5} \emph{Why:} the pilot showed that adding 2.03~ms on the receiver's path lengthened the trip by 1.80~ms for Kafka and by 2.51~ms for Redis, repeatably. The trip is what the law is about, so the delay needed to reach a planned trip has to be measured, not assumed.
\begin{itemize}
  \item \textbf{Steps:} added delays of 0, 1, 2, 4 and 8~ms, plus 16~ms in sessions whose trips reach that far and 32~ms in HZ=100 sessions. The zero step runs twice per round, to catch drift. Each step's delay is measured by the broker's own capture of that run's pings --- how much longer it held the replies to the receiver than the replies to the host --- and that measured value is used.\chgold{D7-1} Ping is recorded beside it and decides nothing: it reads about half a millisecond high, and prints only tenths of a millisecond above 10~ms.\chgold{D6-13} At the 32~ms step, reads come back less often than messages arrive, so trips there do not grow one-for-one for either client.\chgold{D6-6}
  \item \textbf{Order and repeats:} each backend and each load the session runs, 2 rounds before the session's main runs, shuffled as in Section~\ref{sec:stats}, and 2 more when the calibration is too loosely known and nothing else is wrong.\chgold{D6-2}
  \item \textbf{First, a staircase campaign \gl{S0-1}:}\chgold{D4-12} before any other campaign, the steps up to 16~ms run for both backends at 75\% load for 4 rounds, to test with more power whether the line is straight.
  \item \textbf{What it gives:} for each backend, a line of run median trip against measured delay. Each main run's delay is read off that line for its target trip.
  \item \textbf{Checked all session long:} every main run also records its delay and its trip. At analysis the line is refitted from all of the session's runs, and its slope is compared across slices and loads.
  \item \textbf{Gate:} a session's main runs start only if P5(b) holds. Otherwise the session ends, and the pair starts a new one with a new C0.\chgold{D6-11}
\end{itemize}

\textbf{M0. Why the trip does not grow one-for-one \gl{campaign S0-2}.}\chgold{D3-6} \emph{Why:} a law about timestamps needs its measuring setup understood, not just corrected. The departure appears when messages come in bursts \gl{the pilot's replay}, not on the campaigns' steady traffic.\chgold{D6-7} The four candidate explanations, and the test that decides each, are M-H1 to M-H4 in Section~\ref{sec:predictions}.
\begin{itemize}
  \item \textbf{Setups:} the football match replayed in bursts, as in the pilot; both backends, added delays of 0, 2 and 8~ms, at 75\% load on the first x86 pair, 10 rounds; Redis both with one acknowledgement per message and in batches of 200, an intervention on M-H2.\chgold{D6-7}
  \item \textbf{Recording, in half the runs:} a packet capture on the receiver's own address and on the broker; the times at which our receiving program issues each fetch or read and at which it returns; and the receiving thread's wake-up delays \gl{with the tracing tool the harness already uses for the ``got it'' helper}. The other half runs without recording, so the recording's own effect is measured, as in A6.
  \item \textbf{Afterwards:} the request-timing log stays on in later runs only if the recorded and unrecorded halves differ by less than 0.02~ms in median trip and median ``got it'' delay.
\end{itemize}

\textbf{Trip lengths are placed around the predicted cliff.} For each slice $s$ and tick $h$ we use 8 trips: two on the plateau ($0.5s$ and $0.9s$), four across the cliff ($s + 0.2h$, $s + 0.4h$, $s + 0.6h$, $s + 0.8h$), and two beyond it ($s + 1.5h$ and $2(s + h)$). Each trip is reached with the delay the session's calibration says it needs, and the analysis uses the trip each run actually had.\chgold{D3-5}

\textbf{A1. Dose--response on the slice (the main test).} On the first x86 pair, set the time slice by hand to six values: 0.75, 1.5, 2.25, 3, 4.5 and 6~ms, at 75\% load, less those a client cannot reach.\chg{D8-1} \emph{Question:} does the cliff move with the slice, one for one? On the first pair, measured on 18 September, Kafka reaches 2.25, 3, 4.5 and 6~ms and Redis also 1.5~ms; the rest are reported as out of reach with the trip that excludes them.
\begin{itemize}
  \item \textbf{Six campaigns, three per backend:}\chgold{D4-3} slices of 0.75, 3 and 4.5~ms; of 1.5, 3 and 6~ms; and of 2.25 and 3~ms; 8 trips each.
  \item \textbf{The anchor:} the 3~ms slice is in every campaign. Compared across campaigns, its halfway point and plateau height measure each campaign's offset, which the analysis removes before slices are compared \gl{Section~\ref{sec:stats}}.
  \item All six run on one machine pair, in an order drawn with a recorded seed.
\end{itemize}

\textbf{A2. The tick, with three kernels we build ourselves.}\chgold{D3-7} The tick is fixed when the kernel is compiled, so it cannot be changed on a running machine. \emph{Question:} does the cliff's width follow the tick, in proportion?
\begin{itemize}
  \item \textbf{One source, three builds.} The driver's own kernel \gl{6.8.0-1064-azure}, from Ubuntu's source package and with the running kernel's own configuration, is built three times on the same machine with the same compiler, changing only the tick: HZ=1000, 250 and 100 \gl{ticks of 1, 4 and 10~ms}. Our HZ=1000 build is the control, so the three kernels differ in nothing else. Three levels test a proportion; two could only show ``wider''.
  \item \textbf{A bridge to everything else.} Short sessions on Azure's stock kernel run the same setups as our HZ=1000 build, so the tick results connect to A1, A3 and A5, which run on the stock kernel. They should agree within 0.1~ms at the cliff's halfway point.
  \item \textbf{Only the driver changes.} The broker keeps its kernel throughout, because the helper that writes ``got it'' runs on the driver.
  \item \textbf{Checked at every boot, before any run} \gl{manipulation checks}: the tick setting read back from the running kernel; the tick measured by counting the kernel's tick timer on a busy CPU for 10~s, within 5\% of the setting; the scheduler's high-resolution slice timer \gl{HRTICK} switched off in every kernel, because with it on the tick would not set the width and the test would be void; the slice read back as set; and the same tickless settings \gl{NO\_HZ} in all three builds.
  \item \textbf{Booting safely.} A new kernel is booted once while the stock kernel stays the default, so a failed boot falls back to it on the next restart. Azure's serial console is the last resort.
  \item \textbf{Campaigns.}\chgold{D4-7} A build campaign compiles the three kernels and boots each once with its checks. Then twelve sessions, one kernel and one backend each, over four days: each kernel runs with each backend on two different days, and across the days each kernel runs first, middle and last \gl{a Latin square}. Two bridge sessions, one per backend, run on the stock kernel. Each session starts with its own C0 and a short B0 at 75\% load, because the tick may change both the delay's effect and the baseline trip. All fifteen campaigns run on one machine pair.
  \item \textbf{Setups:} 2 slices \gl{1.5 and 3~ms} $\times$ 8 trips placed with that kernel's tick, one backend per session. Each session runs half of the block's rounds, so each kernel and backend gets the full number over its two days.
\end{itemize}

\textbf{A3. Load.}\chgold{D4-4} Run at 50\%, 75\% and 88\% load, at the 3~ms slice and 6 trips, in two campaigns, one per backend. Each campaign calibrates at all three loads. Its rounds come from P3a.\chg{D9-1} \emph{Questions:} does load leave the cliff where it is \gl{P3a}, and does it raise the plateau \gl{P3b}?

\textbf{A4. Another kind of machine.}\chgold{D3-2}\chgold{D4-8} On Azure's Arm machine \gl{Azure Cobalt 100}: first an instrument campaign on the Arm pair \gl{C0, B0 and P0}, then four campaigns, two per backend, at slices of 1.5 and 3~ms and of 3 and 4.5~ms, 8 trips each, less those that pair's own calibration puts out of reach.\chg{D8-1} The 3~ms slice is the anchor, as in A1. \emph{Question:} does the cliff follow the slice on a processor the law never saw?

\textbf{A5. Default slice from core count (secondary).}\chgold{D3-3}\chgold{D4-6}\chg{D11-1} On an x86 driver with 2, 4 and 8 CPUs switched on, left at their default slices, check that the kernel's default follows its rule and that the cliff follows it. The rule is computed from the per-CPU constant each machine reports, not from its version: the first Azure driver's 6.8 kernel carries the constant Linux 6.15 adopted, which gives 1.4, 2.1 and 2.8~ms where the 6.8 rule gives 1.5, 2.25 and 3~ms. Two campaigns, one per backend. Each core count is its own session, with its own C0, at 8 trips.

\textbf{A6. Predict from a recording, with nothing fitted.} In a random half of A1 and A3 runs, record the helper's waits in 0.1~ms steps, and predict the negative rate from that recording. Tracing slows things slightly, so traced and untraced runs are compared too.

\textbf{A7. Go-first priority.}\chgold{D4-5} At the 3~ms slice and three trips \gl{$0.9s$, $s + 0.5h$ and $s + 1.5h$}, turn go-first priority on for the timestamping helpers only, at the same load. Both backends, in one campaign.

\textbf{A8. Client language.}\chgold{D4-9} Kafka only: our Python client against Kafka's official Java client, with the ``got it'' note taken on the background helper or on the sending thread, at the 3~ms slice and the same three trips as A7, ordinary and go-first. One campaign on an x86 pair and one on the Arm pair. Each client has its own calibration, because the delay's effect on the trip belongs to the client.

\section{Experiments for Mode 2 (the industry)}

Each runnable tool is one campaign on one machine pair, with our own program recording the same messages.\chgold{D4-1} For each tool, at least two experiments, and three where it applies:

\textbf{T1. Delay staircase} \gl{receiver-only, as in the fix above}. Add 0, 0.1, 0.2, 0.3, 0.5, 0.7, 0.9, 1.1, 1.5 and 2.0~ms, in random order. Record the tool's average, p50 and p99, and our reference at each step. \emph{Questions:} does the tool see a 0.1~ms step? Does it round, chop or keep? Does its average track the truth?

\textbf{T2. Forced negatives.} Make the receiver's clock read 0.5~ms and 2~ms early \gl{libfaketime for Java, C and Python tools, checked on each tool first; a second machine with a deliberately offset clock for Go tools, restored afterwards}. \emph{Questions:} does it average over the negatives, drop them, replace them with zero, crash, or count them? Does its message count still match what was sent?

\textbf{T3. Load and go-first priority.} Run each tool on an idle machine and at 88\% load, with and without go-first priority on the tool's own process. \emph{Question:} how much of the tool's own number is waiting in line?

\textbf{T4. What it measures.} Check which trips each tool can report at all: send to receive, send to ``got it'', or only request to reply. Confirm from the documentation, the source and a real output file.

\textbf{T5. Documentation and issue audit} \gl{desk work, no machines}.\chgold{D3-8} For each tool in Section~\ref{sec:tools}:
\begin{enumerate}
  \item Read its documentation for three things: which clock and resolution it times with; what it does with zero or negative values; and any warning about scheduling or clock effects.
  \item Search its issue tracker and mailing list with fixed terms: ``negative latency'', ``negative'', ``clock skew'', ``resolution'', ``rounding'', ``millisecond'', ``nanosecond'' and ``latency 0''.
  \item Record every relevant report with its date and outcome: fixed, closed without a fix, or still open.
\end{enumerate}
The terms and the reading rules are frozen with this plan. Results are counts per tool, with every report linked. \emph{Question:} where a tool mishandles values, does anyone say so, and did anything change when users noticed?

\section{What is already known}
\textbf{Mode 1: no law found.}
\begin{itemize}
  \item \textbf{The kernel's rule is only a bound.} The developers state a bound, slice plus tick, not a distribution. Their 2026 conference talk measured wakeup delays and explicitly presented no model.
  \item \textbf{Classic waiting-time laws cover jobs, not timestamps.} Kleinrock (1967) and Coffman, Muntz and Trotter (1970) give waiting-time laws under time slices for jobs in a queue; nobody has applied them to timestamps.
  \item \textbf{Nearby laws cover other things.} Tsafrir (2005) models noise from kernel ticks in parallel programs. Sharma (2026) gives a condition for clock differences, not for scheduling.
  \item \textbf{Timestamping and tracing studies are empirical.} Villain (2012), Giraldeau and Dagenais (2015) and Bielmeier (2025) measure delays or predict them from recordings, with no formula from settings.
\end{itemize}

\textbf{Mode 2: comparisons exist outside brokers.} OctoPerf (2022, practitioner blog) found four HTTP tools reporting averages from 20 to 150~ms on one site. Lancet (2019) found three load testers reporting 99th percentiles of 27, 40 and 63~µs. Every broker comparison found uses one tool per study, and none checks rounding, negatives or discards.

\textbf{Practitioner sources (ordinary web search).}
\begin{itemize}
  \item \textbf{Why tools disagree.} Practitioner write-ups agree that tools report different latencies mainly because they start and stop the clock in different places. OctoPerf's comparison of four HTTP tools is the clearest case. The ``coordinated omission'' write-ups make the same point about \emph{when} the clock starts: Gil Tene's wrk2, Tyler Treat's 2016 message-queue benchmarks, ScyllaDB (2021) and the YCSB fix. Kroxylicious (2026) notes that Kafka's producer perf test cannot show end-to-end latency at all.
  \item \textbf{Millisecond clocks.} It is widely known that a millisecond clock cannot time sub-millisecond work \gl{Shipilev's ``Nanotrusting the Nanotime'', the LMAX blog, many Java write-ups}. None of these connects it to results being silently thrown away inside a benchmark's percentiles.
  \item \textbf{Negative latency in practice.} Negative latencies and negative span durations are reported often: a Kafka mailing-list thread, Jaeger and OpenTelemetry users, and a 2026 Kafka monitoring guide. They are almost always blamed on clock skew; RabbitMQ PerfTest's own documentation warns about clock sync only. Tools deal with them by reporting ``not a number'' \gl{Kafka's consumer latency metric}, clamping to zero \gl{Splunk's OpenTelemetry collector}, dropping them, or breaking dashboards.
  \item \textbf{Scheduler delays.} Kernel developers and practitioners report latency spikes and regressions tied to the EEVDF scheduler's slices \gl{LWN 2026; a 2026 write-up on fixing EEVDF regressions at Meta}. None connects those delays to wrong timestamps in latency tools, and none gives a law.
  \item \textbf{Bottom line.} Practitioners know the symptoms. None of the sources we found ties them to scheduling or to a predictive law, and none compares broker latency tools against each other. T5 makes this a count rather than an impression.\chgold{D3-8}
\end{itemize}

\section{Predictions, frozen before running}\label{sec:predictions}
{\small
\begin{longtable}{@{}p{1.0cm} p{9.5cm} p{5.6cm}@{}}
\toprule
\textbf{ID} & \textbf{Prediction} & \textbf{Counts as confirmed if} \\
\midrule
\endhead
P1 & \textbf{The cliff follows the slice.} For each slice $s$ the pair can reach,\chg{D8-1} the trip at which the negative rate falls halfway from the plateau \gl{the level at $0.9s$} to the floor \gl{the level at $2(s+h)$} lies between $s$ and $s + h$. Across those slices, after each campaign's offset is removed through the anchor slice,\chgold{D4-3} that halfway point rises one-for-one with $s$. & In all but one of the slices the pair can reach, and at least four such slices, the halfway point is inside the band.\chg{D8-1} The straight-line slope against $s$ lies between 0.8 and 1.2, with a 95\% interval that excludes 0.5 and 1.5. A slice out of reach is reported with the zero-delay trip that excludes it. \\
P2 & \textbf{The cliff's width follows the tick.} The drop from 90\% to 10\% of the plateau is about four times wider at HZ=250 than at HZ=1000. & The width ratio's 95\% interval overlaps 2.5--5.5 and excludes 1. \\
P2b & \textbf{And about ten times wider at HZ=100.}\chgold{D3-7} & The width ratio's 95\% interval overlaps 6--14 and excludes 1. \\
P2c & \textbf{The width is proportional to the tick.}\chgold{D3-7} Across the three kernels, the fitted cliff width against tick length is a straight line through zero with slope one. & The slope's 95\% interval overlaps 0.8--1.2 and excludes 0.5; the intercept's 95\% interval includes 0. \\
P2d & \textbf{The tick does not move the cliff's start.}\chgold{D3-7} The fitted start stays at the slice whatever the tick. & In all three kernels, at both slices, the fitted start is within 0.25~ms of the slice. \\
P3a & \textbf{Load leaves the cliff where it is.}\chgold{D9-1} At the 3~ms slice,\chgold{D4-4} the halfway point moves by less than 0.25~ms between the lowest and the highest load whose ordinary plateau is at least 2\%.\chg{D14-2} & The whole 95\% interval of the movement lies inside $\pm$0.25~ms \gl{two one-sided tests, on the signed difference}. A point estimate inside the band is not enough: a scattered estimate can land inside it by luck, and this is a claim that something did not happen. A load below 2\% is reported as out of reach, with the plateau that excludes it, and is not judged; where fewer than two loads remain, P3a is reported as out of reach on that pair.\chg{D14-3} \\
P3b & \textbf{Load raises the plateau.}\chg{D9-1} At the same slice and between the same loads, the plateau height rises. & The increase has a 95\% interval above zero. Judged on its own; it is reported as underpowered where the campaign's rounds cannot reach the power the rule asks of it. \\
P4 & \textbf{Go-first priority removes the plateau.} At the 3~ms slice,\chgold{D4-5} it cuts the plateau rate at least 5 times wherever the ordinary plateau is at least 2\%. & Point estimate $\ge 5$, 95\% interval above 2. \\
P5 & \textbf{Measurement checks.}\chgold{D3-5} (a) Arrival minus sending is never below zero. (b) The receiver-only delay lengthens the trip in a way we can rely on. (c) It leaves the ``got it'' delay unchanged. & (a) 0 negatives per run. (b) In every session's calibration, the slope's 95\% interval lies above 0.5, and at every step the calibration's 95\% interval lies within 0.3~ms of it \gl{from the scatter of the runs about the line, or within the steps where the step medians are joined}; each run's and each step median's distance from it is reported.\chgold{D6-1} (c) Equivalent within $\pm$0.10~ms on the median and $\pm$10\% on the p99, at every delay. Otherwise the session is analysed with each run's own ``got it'' distribution, and the departure is reported. A median shift stops the session when it is larger than a quarter of the added delay \emph{and} larger than the session's own got-it noise --- 0.10~ms, or three times the scatter of its zero-delay runs, whichever is larger.\chg{D8-2} \\
P6 & \textbf{A recording predicts the rate}, with nothing fitted, within a factor of 1.5, and not systematically too low. & At least two thirds of setups within 0.67--1.5, and the median ratio at least 0.8. \\
P7 & \textbf{Default slices follow the core-count rule}, computed from the per-CPU constant read off each machine \gl{on the first Azure machine: 1.4, 2.1 and 2.8~ms at 2, 4 and 8 cores}, and the cliff follows them.\chgold{D3-3} & Read-back values match; P1's band holds at each. \\
P8 & \textbf{Language (Kafka).}\chgold{D4-9} Java shows the same cliff and priority effect as Python: the rate at $s + 0.5h$ lies between its levels at $0.9s$ and $s + 1.5h$, and go-first cuts the plateau at least 5 times. Python's rate is at most 1.5 times Java's at matched setups. & All hold at the 95\% level, on each machine pair. \\
P9 & \textbf{New hardware.}\chgold{D4-14} On the Arm pair the cliff follows the slice: at each slice that pair can reach\chg{D8-1} the halfway point lies between $s$ and $s + h$, and it rises with the slice, after the anchor correction. & The halfway point is inside the band at every slice the pair can reach, and at least two such slices, and its straight-line slope against $s$ lies between 0.8 and 1.2. \\
\midrule
\multicolumn{3}{@{}p{16.1cm}@{}}{\emph{Why the trip does not grow one-for-one (M0).\chgold{D3-6} Each explanation is kept or dropped by its test, and more than one may hold. ``The departure'' is how far the measured slope is from one, times the added delay.}} \\
M-H1 & \textbf{The network.} The delay that reaches the receiver's network card differs from the delay set. & Kept if the arrival shift in the packet capture differs from the set delay by more than 0.05~ms. \\
M-H2 & \textbf{The receive loop.} A longer delay changes whether a new message finds a request already waiting at the broker, or has to wait for the next one. & Kept if a replay of each run's message times against its logged request cycle reproduces the measured slope within 0.05. \\
M-H3 & \textbf{The wake-up.} With a delay, the receiving thread has gone to sleep when the reply arrives, and must be woken and scheduled. & Kept if the change in the thread's median wake-up delay between 0 and 8~ms accounts for at least half of the departure. \\
M-H4 & \textbf{The TCP acknowledgement} \gl{expected for Redis}. Delaying the broker's replies also delays its acknowledgements, which can hold the receiver's next request back. & Kept if requests held for an acknowledgement appear in the capture and account for at least half of the departure. \\
\midrule
T-P1 & \textbf{Low, chops after measuring} (Kafka EndToEndLatency). Percentiles equal the true value rounded down to whole ms: 0 for every step below 1~ms. The average is within 2\% of our reference. & All staircase steps match. \\
T-P2 & \textbf{Low, reads a ms clock twice, keeps all} (ProducerPerformance, emqtt-bench). With send times spread evenly, the average is within 0.05~ms of the truth, and the median flips from 0 to 1 at a true 0.5~ms. With sends paced in step with the tick, the average can be off by up to 0.5~ms. & Both patterns seen at the predicted steps. \\
T-P3 & \textbf{Low, drops $\le 0$} (OpenMessaging end-to-end). For trips under 1~ms, the median prints 1.0~ms at every step, and retention is about the trip divided by 1~ms. & Median constant below 1~ms; retention within 5 points of prediction. \\
T-P4 & \textbf{Low, drops negatives, keeps zeros} (Pulsar perf). The median is 0 below a true 0.5~ms. Under a forced offset, the average rises by exactly what dropping the negatives removes. & Average shift within 10\% of the computed value. \\
T-P5 & \textbf{Medium} (wrk2). It tracks 0.1~ms steps within 0.05~ms at the median. Forced negatives are not recorded, or the run aborts. & Both behaviours as predicted. \\
T-P6 & \textbf{High, keeps all} (RabbitMQ PerfTest, NATS CLI). They track 0.1~ms steps within 0.02~ms. A forced offset shifts every value by the offset, and negatives show in the minimum and pull the average down. & Shift within 0.02~ms; negatives present. \\
T-P7 & \textbf{Every ``read first'' tool} behaves as its class predicts, in T1 and T2. & At least 90\% of the predicted behaviours observed. \\
T-P8 & \textbf{The documentation stays silent.}\chgold{D3-8} Where the source rounds, drops, clamps or replaces values \gl{T-P1 to T-P5, fio, btt}, the documentation does not say so. & Documented in at most a third of those tools. \\
\bottomrule
\end{longtable}}

\section{Statistics: how many runs, and in what order}\label{sec:stats}

\textbf{How many rounds, from each campaign's own test.}\chgold{D4-2} The number of rounds is not fixed in advance; the rule that sets it is. Before a campaign's main runs, and after its machine pair's spread pilot \gl{P0}, the campaign is simulated 1,000 times on made-up data:
\begin{itemize}
  \item \textbf{The made-up data follow the law exactly.} They use the plateau height and floor that the pair's P0 and B0 measured at the campaign's load; the campaign's own slices, trips, loads and tick; the planned number of messages per run; and run-to-run variation, on the logarithm of the rate, equal to the spread P0 measured on that pair for that backend --- one figure per backend, never the two pooled, because pooling assumes a shared variance that the pilots refute.\chg{D10-1} Shifts between campaigns and between sessions are added, with a standard deviation of 0.1~ms in trip and 20\% in plateau height.
  \item \textbf{Each simulated campaign goes through the real analysis} and the same ``counts as confirmed if'' rule as the real data (Section~\ref{sec:predictions}).
  \item \textbf{Rounds} are the smallest number, at least 4, for which the rule confirms the prediction in at least 80\% of the simulations, and in at most 5\% of simulations where the prediction is false in the way its falsifier states: a cliff that stays put as the slice changes \gl{P1, P9}; a width that does not change with the tick \gl{P2 to P2c}; a halfway point that moves 0.5~ms between 50 and 88\% load \gl{P3a}, judged over the loads that remain\chg{D14-2}; no go-first effect \gl{P4}; a cliff that stays put as CPUs are switched off \gl{P7}; Python's rate twice Java's \gl{P8}.
  \item \textbf{Blocks of several campaigns} \gl{A1, A2, A4} are simulated as a whole, and every campaign of a block runs the same number of rounds.
  \item \textbf{Ceiling:} 40 rounds. If 40 rounds cannot reach that power, the campaign runs at 40 and its prediction is reported as underpowered, as written in the campaign's log before it runs.
  \item \textbf{Fixed numbers where no prediction is tested:} C0 runs 2 rounds per session and 4 in S0-1; B0 runs 3; P0 runs 5; M0 runs 10. Tool campaigns run 4 rounds: their predictions are patterns \gl{rounding, dropping, a median stuck at 1~ms} that one round already shows, and the repeats measure how stable each pattern is.
  \item \textbf{Why not a fixed floor.} Version 3 required at least 15 runs per setup, from old setups whose run-to-run spread was about as large as the rate. The pilot's repeat runs varied far less \gl{a standard deviation of 0.14 to 0.23 on the log rate}, and a simulation of the actual decision rule measures power directly. The simulation code is tested on made-up data before its first use, and its seed is recorded.
\end{itemize}

\textbf{Spread pilot (P0):} 5 rounds of 16 setups at 75\% load on every machine pair, before that pair's main campaigns.\chgold{D4-10} They are used only to estimate run-to-run spread, never to test a prediction. The instrument pilot in Section~\ref{sec:pilot} is a different thing.

\textbf{Machine pairs.}\chgold{D4-10}
\begin{enumerate}
  \item \textbf{One pair per test.} Every prediction is tested on runs from one pair. A1, A3 and A7 run on the first x86 pair; A5 and all fifteen A2 campaigns on one x86 pair; A4 on the Arm pair; each tool on one pair. A8 runs once on an x86 pair and once on the Arm pair, and P8 is judged on each.
  \item \textbf{Every pair opens with its own B0 and P0} and reads back its settings, so the pairs can be described against each other and the repeats rule uses that pair's own spread.
  \item \textbf{Every run records its pair.} No estimate pools runs from different pairs.
  \item \textbf{At the same time:} campaigns on different pairs may run at once; campaigns on one pair never overlap.
\end{enumerate}

\textbf{Random order, so a passing server mood cannot fake an effect:}
\begin{enumerate}
  \item \textbf{Rounds.} List every setup of the campaign. One round contains every setup once. Repeat for as many rounds as the rule gives.
  \item \textbf{Shuffle.} Within each round, shuffle the order with a random generator whose seed is written in the run log.
  \item \textbf{No back-to-back repeats.} If the last setup of one round equals the first of the next, reshuffle.
  \item \textbf{Spread over time.}\chgold{D4-1} The campaigns of one block run in an order drawn with a recorded seed, over at least two days and at different times of day \gl{morning, afternoon, night}.
  \item \textbf{Settings that need a reboot, a new machine or another pair.} Tick, kernel, core count, machine type and machine pair cannot be shuffled run by run. They are \emph{sessions}: the order of sessions is randomized, and each session type is repeated on at least 2 different days. Kernel sessions are counterbalanced across days.\chgold{D3-7} The slice \emph{can} be changed without a reboot, so it is shuffled run by run inside a session like everything else. The analysis treats session as its own source of variation \gl{split-plot design}.
  \item \textbf{Failed runs.}\chgold{D5-2} If a run fails mechanically, or a condition it was meant to have did not take \gl{Section~\ref{sec:operations}}, it is logged with the reasons and re-queued at a random later position in the same session, 3 attempts in all. A run whose conditions held but that gives a surprising number is a result, never a re-run.
\end{enumerate}

\textbf{Integrity checks, the same for every arm.}\chgold{D5-2} Before a run counts, it must pass all of these. Section~\ref{sec:operations} says when each is checked and what follows a failure. Runs that fail are counted, reported and never silently dropped.
\begin{itemize}
  \item Arrival minus sending is never negative.
  \item The session's delay calibration passed P5(b), and the run shows no gross ``got it'' departure \gl{P5(c)}.\chgold{D3-5}
  \item At least 99\% of the planned messages are sent after the warm-up, and at least 99\% of those arrive.\chgold{D5-2}
  \item The achieved send rate is within 2\% of target.
  \item Measured load is within 3 points of target, from at least 10 samples taken while messages are sent.
  \item The slice and tick read back as intended, before and after the run.
  \item The added delay, from the broker's own capture of the run's pings, is within 0.05~ms of the delay set, and that capture exists. What ping made of it is recorded.\chgold{D5-2}\chgold{D6-5}\chgold{D7-1}
  \item Before the run, the machine is as its session set it up: the receiver's address only in its namespace, the broker answering, and no package manager running.\chgold{D6-8}
  \item The clients ran with the campaign's settings, as their own logs state: Kafka's producer with up to 64 requests in flight, Redis's consumer acknowledging in batches of 200.\chgold{D6-6}
  \item The clock's offset is logged before and after the run.
  \item The first 30 seconds are discarded.
\end{itemize}

\textbf{How the cliff is measured, fixed before the data.} For each slice, the curve of negative rate against trip is summarized two ways:
\begin{itemize}
  \item \textbf{Fitted shape:} fit the plateau, straight fall and floor to the runs, and read off where the fall starts and how wide it is.
  \item \textbf{Free shape:} fit a smooth curve that is only required to go down, and read off the halfway point and the 90\%-to-10\% width.
\end{itemize}
The law passes only if both agree with the prediction. Intervals come from resampling whole rounds \gl{bootstrap}. \textbf{Each run sits at the trip it actually had}, its own measured median trip, not the trip it was aimed at.\chgold{D3-5} As a secondary analysis, fixed now, every message is also placed at its own trip within its session, which traces the cliff inside runs as well as across them.

\textbf{Removing campaign offsets through the anchor.}\chgold{D4-3} In A1 and A4, a campaign's offset is its anchor slice's halfway point minus the mean of the anchor's halfway points over the block's campaigns; for plateau height, the same on the log scale. Each campaign's other halfway points and plateau heights are corrected by its offset before P1 or P9 is judged. Both predictions are also reported without the correction.

\textbf{Analysis, fixed before the data.}
\begin{itemize}
  \item \textbf{Unit of analysis:} the run.
  \item \textbf{Comparisons:} ratios on the log scale, allowing for session and round.
  \item \textbf{Pairs and campaigns:}\chgold{D4-10} machine pair, campaign and session are sources of variation; nothing pools runs across pairs.
  \item \textbf{Several predictions in a family:} tested together with Holm's correction.
  \item \textbf{Calibration} \gl{C0}:\chgold{D3-5} for each backend and session, a straight line of run median trip against measured added delay, with a test of the line against the per-step medians \gl{lack of fit}. Intervals come from resampling whole rounds. If the line fails the test, the per-step medians are joined by straight segments instead, and the shape is reported.
  \item \textbf{``Got it'' equivalence} \gl{P5(c)}:\chgold{D3-5} two one-sided tests at 90\% at each delay step against zero delay, with margins of $\pm$0.10~ms on the median and $\pm$10\% on the p99.
  \item \textbf{Code first:} the analysis code is written and tested on made-up data before the first real run it analyses.
  \item \textbf{Frozen predictions:} this plan is frozen as Section~\ref{sec:freeze} describes, before the first campaign.
  \item \textbf{Every run tracked:} failures included, as in all our earlier campaigns, each with its verdict, its reasons and its fingerprints.\chgold{D5-4}
\end{itemize}

\clearpage
\section[The campaigns]{The campaigns\protect\chgold{D4-1}}\label{sec:campaigns}
{\small
\begin{longtable}{@{}>{\raggedright\arraybackslash}p{2.2cm} >{\raggedright\arraybackslash}p{5.0cm} >{\raggedright\arraybackslash}p{3.7cm} >{\raggedleft\arraybackslash}p{1.1cm} >{\raggedleft\arraybackslash}p{1.5cm} >{\raggedleft\arraybackslash}p{1.1cm}@{}}
\toprule
\textbf{Stage} & \textbf{Campaigns} & \textbf{Question} & \textbf{Runs} & \textbf{Hours each} & \textbf{Hours} \\
\midrule
\endhead
0 Instrument & 4: S0-1 delay staircase; S0-2 M0; S0-3 B0 and P0; S0-4 B0 and P0 on a second x86 pair, if one is used & how the delay moves the trip, why, and how much runs vary & 414 & 3.5--7.4 & 25.2 \\
1 Slice \gl{A1} & 6: per backend, slices of 0.75, 3 and 4.5~ms; 1.5, 3 and 6~ms; 2.25 and 3~ms & does the cliff follow the slice? & 592 & 4.7--6.7 & 36.0 \\
2 Load \gl{A3} & 2: one per backend & does load leave the cliff where it is \gl{P3a}, and raise the plateau \gl{P3b}? & 576\chg{D9-3} & 6.5 & 33.6 \\
3 Go-first \gl{A7} & 1 & does priority remove the plateau at the cliff? & 72 & 4.5 & 4.5 \\
4 Core count \gl{A5} & 2: one per backend & does the kernel's default slice set the cliff? & 264 & 8.4 & 16.9 \\
5 Tick \gl{A2} & 15: the build, 12 sessions, 2 bridge sessions & does the cliff's width follow the tick? & 682 & 3.5--5.2 & 55.5 \\
6 Arm \gl{A4} & 5: the Arm instrument campaign, then 2 per backend & does it hold on another processor? & 430 & 4.7--7.4 & 26.3 \\
7 Client \gl{A8} & 2: an x86 pair and the Arm pair & is it the Python client? & 288\chg{D13-1} & 8.6 & 17.3 \\
8 Tools \gl{T1--T4} & 10: one per runnable tool & what do the tools report? & 680 & 4.7 & 47.2 \\
T5 audit & desk work & does anyone say so? & -- & -- & -- \\
\midrule
\textbf{Total} & \textbf{47} & & \textbf{3,638}\chg{D13-1} & & \textbf{242.0} \\
\bottomrule
\end{longtable}}

Runs and hours are at the floor of 4 rounds, except where a campaign's own rounds have been
simulated and the row carries the change that set them: A3 at 16 rounds \gl{D9-3} and A8 at 6
\gl{D13-1}. A run takes 3.5~minutes, as the pilot measured \gl{warm-up, at least 5,000 messages, settling and checks}; each campaign adds 0.25~hours of setup, and each reboot 0.5~hours. At 6 rounds the total is 4,826 runs and 312~hours. M0's second Redis arm adds 30 runs \gl{D6-7}; a calibration's second stage, where one is needed, adds 24 runs to its session \gl{D6-2}.

\textbf{Which pair runs what.} With one x86 pair, every x86 campaign runs on it and S0-4 is not needed. With two x86 pairs, the first runs S0-1 to S0-3, A1, A3, A7, A8's x86 campaign and five tool campaigns; the second runs S0-4, A5, all of A2, and the other five tool campaigns. The Arm pair runs A4 and A8's Arm campaign.

\textbf{How long.} At about 20 hours a day: one x86 pair and then the Arm pair, about 11.5 days; two x86 pairs at once and then the Arm pair, about 7 days; all three pairs at once, about 5 days.\chgold{D5-7} Each new pair's shakedown adds a few hours to these. The rules of Sections~\ref{sec:stats} and~\ref{sec:operations} hold in every case.\chgold{D5-1}

\textbf{Money.} About EUR 97 of machine time at list prices with all three pairs in Sweden Central \gl{EUR 126 at 6 rounds}, plus about EUR 26 a month per pair for disks and fixed addresses while the pair exists.\chgold{D5-7}

\textbf{Order of work:}
\begin{enumerate}
  \item \textbf{Now, with the machines stopped.} Freeze this plan, together with the code that judges each run, stops a campaign, watches the pairs and copies the data, which is written and tested before the freeze.\chgold{D5-2} Write, and test on made-up data, the code for M0's recording, the tick checks, the simulated-power rule and the law analysis. Read the ``read first'' tools and freeze their classes. Do the T5 audit. Prepare the kernel builds and the Java client.
  \item \textbf{Stage 0.} On the first x86 pair: S0-1, then S0-3, then S0-2 once M0's recording is tested. On a second x86 pair, if one is used: its shakedown, then S0-4 and A2's build. The Arm pair's shakedown comes before its instrument campaign.\chgold{D5-1} Stage 0 starts again from the beginning under version 6. After any change to the code, the first x86 pair runs alone until its shakedown and first calibration pass, and then the others start.\chgold{D6-12}
  \item \textbf{Rounds.} Each campaign's rounds come from Section~\ref{sec:stats} and are written in its log before it runs.
  \item \textbf{The campaigns,} in the order of the table above on each pair, with each block's campaigns in their seeded order. Each campaign's runs are copied and checked before the next campaign starts on its pair.\chgold{D5-4}
  \item \textbf{Analysis and report.} Every prediction reported as confirmed or not, with every run listed.
\end{enumerate}
\textbf{If money or time runs short:} the slice and tick campaigns first, since they are the law; then A4 and the tools; then the rest.

\section[How the campaigns are run, checked and stopped]{How the campaigns are run, checked and stopped\protect\chgold{D5-1}}\label{sec:operations}
Campaigns run unattended, often at night, on up to three machine pairs at once. So the rules for running them are fixed here, before any run, and a program applies them, not a person looking at results.

\textbf{Lanes.}\chgold{D5-6} Each machine pair is a \emph{lane} \gl{a driver and its broker, with their own queues and logs}. A campaign runs on one lane from its first run to its last, and the campaigns of one lane never overlap. Two lanes share no machine, broker, network link or disk, so a run on one cannot disturb a run on the other. Every lane runs the same version of the code, and every run records its lane, its driver and that version \gl{the commit} next to its files.

\textbf{A new pair's shakedown.}\chgold{D5-1} Before a pair's first campaign, the pilot of Section~\ref{sec:pilot} runs on it. Its scheduler-settings, receiver-only-delay, load and never-negative checks must pass, and a pair that fails one runs nothing until the cause is found. The receiver-only delay is checked by the broker's own capture of the pilot's pings, within 0.05~ms;\chgold{D6-3} with no delay, the two paths are measured by ping, TCP and UDP and recorded, and both machines' network and timing facts are kept with them.\chgold{D7-2}\chgold{D7-3} A pair that passed its shakedown repeats this network part in every later session, because a start after a stop can put a machine on another physical host.\chgold{D6-4} The pilot's harness-delay check is replaced by the pair's own C0, as on the first pair \gl{D3-5}. Its go-first result is recorded but is not a gate, because A8 tests go-first as a prediction.

\textbf{Every run, as it ends.}\chgold{D5-2} A program on the driver checks each run the moment its messages stop, writes every check next to the run with the value found and the limit it was held to, and gives one of three verdicts:
{\small
\begin{longtable}{@{}>{\raggedright\arraybackslash}p{2.1cm} >{\raggedright\arraybackslash}p{8.3cm} >{\raggedright\arraybackslash}p{5.4cm}@{}}
\toprule
\textbf{Verdict} & \textbf{When} & \textbf{What follows} \\
\midrule
\endhead
Counts & Every integrity check of Section~\ref{sec:stats} passes. & The run enters the analysis. \\
Repeated & A condition the run was meant to have did not take: fewer than 99\% of the planned messages sent after the warm-up, or fewer than 99\% of those arrived; the send rate more than 2\% from target; the load more than 3 points from target, or fewer than 10 load samples while messages were sent; the slice or tick different after the run from what was set; the added delay, from the broker's capture, more than 0.05~ms from the delay set, or that capture missing;\chgold{D7-1} the machine not as its session set it up before the run;\chgold{D6-8} a client not set as every campaign sets it;\chgold{D6-6} the clock's offset not logged before and after; or a step that failed, or files that cannot be read. & The run keeps its files and its line in the ledger, marked failed with its reasons, and a copy is queued at a random later place in the campaign, 3 attempts in all. A run whose third attempt fails is abandoned and reported. \\
Stops the campaign & The instrument is in doubt: a message arrived before it was sent \gl{P5(a)}, or the run's ``got it'' median moved from the session's zero-delay median by more than a quarter of the delay added \gl{P5(c)}. & The run is recorded as failed, and nothing more runs on that pair until the cause is found and recorded with the campaign. If the cause is in the code every pair runs, the other pairs stop too. \\
\bottomrule
\end{longtable}}
Nothing a run measured makes it a repeat: a run whose conditions held is a result, whatever its number. The ``got it'' comparison needs the session's calibration, so it applies to every run after the session's C0; within C0 itself, P5(c) is judged when the calibration is fitted.

\textbf{When attempts keep failing.}\chgold{D5-3} After every attempt the ledger is read again. The campaign stops itself when the last 3 attempts all failed, or when more than 20\% of the last 20 failed, once at least 10 have finished. A machine that has gone wrong then costs minutes, not a night.

\textbf{Void sessions.}\chgold{D6-10} A session broken by a fault of the machines or of our code, or stopped to fix one, before its calibration's gate is void. Nothing in it is deleted: its files and ledger are copied and fingerprinted like any campaign's, a note beside them gives the cause, and the report lists it. It is never analysed. A session whose calibration fails its gate ends too, and the pair starts a new one with a new C0; no session is held open while a rule changes.\chgold{D6-11}

\textbf{Watching every pair.}\chgold{D5-5} A watch on a computer outside the testbed reads every pair every 5 minutes. It changes nothing on the machines, except that it deallocates a pair whose machines have been idle for 20 minutes, so that the pair bills only for its disks. It raises:
\begin{itemize}
  \item \textbf{alerts,} which need a person now: a campaign that stopped itself, or a run judged to stop one; a message that arrived before it was sent; a campaign with no progress for 20 minutes; new failed runs; a broker that is not running; a disk 80\% full, or less than a tenth of memory free; a machine that does not answer; one queue running on two pairs;
  \item \textbf{warnings,} worth a look: a repeated run; CPU time taken by other tenants above 5\% \gl{steal}; a clock more than 1~ms off; load far below its setting; pairs running different versions of the code;
  \item \textbf{notices:} a finished campaign, whose runs are to be copied before the next campaign starts on that pair.
\end{itemize}
Every look is kept in a dated log.

\textbf{Keeping the data.}\chgold{D5-4} A campaign's runs exist only on its driver until they are copied. When a campaign finishes, and before the next starts on that pair, its ledger and every file of every attempt, done or failed, are fingerprinted on the driver \gl{SHA-256: a 64-character code that changes if any byte of the file changes}, packed, copied, and checked file by file against their fingerprints. The fingerprints are kept with the data, with a record of the machine, the code version and the time of the copy. A copy that fails its check is made again from the driver, where nothing is deleted. A pair is deleted only after all its campaigns are copied and checked. What a driver measured outside any campaign \gl{a pilot, a void session, a probe of the instrument, a log} is copied and checked the same way, and before a pair is deleted its whole folder of runs is copied once more. Every run also keeps the TCP counters of the driver, the receiver's namespace and the broker, and the broker's own log for the run.\chgold{D6-10} Every run writes into a folder of its own, named after its campaign and its place in it, and a campaign that finds that folder already there stops.\chgold{D6-9}

\section{What could go wrong, and what it would mean}
\begin{itemize}
  \item \textbf{P1 fails} (the cliff does not follow the slice). The ``law'' is only a description of our old machines. That is still worth reporting, and it would redirect the paper.
  \item \textbf{P2 to P2c fail} (the width does not follow the tick). The wait is not ``slice plus up to one tick''. The plateau--cliff shape may survive with a different width rule. If only P2c fails, the width grows with the tick but not in proportion, and something else sets it at long ticks.
  \item \textbf{P2d fails} (the tick moves the cliff's start). The slice is enforced only on tick boundaries, and the law's start becomes the slice rounded up to a tick.\chgold{D3-7}
  \item \textbf{P3 fails} (load moves the cliff). Waiting is more than one protected slice under heavy load, for example several tasks queued ahead. That would point to the classic time-sharing queueing laws, which do predict this.
  \item \textbf{P8 fails on Java.} The effect is tied to Python's lock or our client. That narrows Mode 1 to Python-style clients.
  \item \textbf{The slice setting does not stick}, because the kernel resets it when cores change. Read it back before every run and discard runs where it drifted.
  \item \textbf{The delay's effect is not a straight line} (C0).\chgold{D3-5} The per-step medians are joined instead, and the shape is reported. If a longer delay can give a shorter trip, the receiver-only delay cannot place trips, and that machine's main runs do not start.
  \item \textbf{No explanation survives M0.}\chgold{D3-6} The slope is reported as an open property of the clients. The law's tests still stand, because each run is analysed at the trip it actually had.
  \item \textbf{HRTICK is on in a kernel.}\chgold{D3-7} It is switched off, or the kernel rebuilt, before any tick run. A tick test with it on would be void.
  \item \textbf{A new kernel does not boot.}\chgold{D3-7} The machine falls back to the stock kernel, the failure is logged, and no run is lost.
  \item \textbf{The two backends disagree.} Both are reported, and the law's claim is narrowed to where it holds.
  \item \textbf{The machine pairs differ}\chgold{D4-10} in baseline trip or run-to-run spread. Each pair's tests stand on their own, with that pair's own repeats, and the difference is reported.
  \item \textbf{A simulation asks for more than 40 rounds.}\chgold{D4-2} The campaign runs at 40, and its prediction is reported as underpowered, as its log says before it runs.
  \item \textbf{The anchor moves between campaigns}\chgold{D4-3} by more than 0.25~ms at its halfway point. The movement is reported, and P1 is given both with and without the anchor correction.
  \item \textbf{A campaign stops itself.}\chgold{D5-3} Nothing more runs on that pair until the cause is found. Runs finished before the stop stand if the cause did not touch them; if it did, they are repeated, and both the cause and the repeats are reported.
  \item \textbf{A pair is lost in the middle of a block}\chgold{D5-4} \gl{a machine fails, or its disk is lost}. The campaigns already copied are safe. Because no test pools runs from two pairs, the interrupted block runs again in full on a rebuilt pair, after its shakedown, B0 and P0, and the interrupted runs are reported.
  \item \textbf{A pair's runs vary too much for its calibration.}\chgold{D6-2} Two more rounds run; if the calibration is still too loosely known, the session ends, is reported with its spread, and the pair starts a new one.
  \item \textbf{The two paths differ with no delay.}\chgold{D7-2} It is reported with the session, in all three measurements, and the pair runs: what its runs vary by is what its rounds are set from. On 17 September a gate on ping's figure turned away a pair whose messages were unaffected.
  \item \textbf{A run has no capture of its delay.}\chgold{D7-1} The run is repeated. If captures keep failing, the campaign stops itself on the run of failures \gl{D5-3}, because a treatment nobody measured is not a treatment we can report.
  \item \textbf{A client does not take its setting.}\chgold{D6-6} Its runs are repeated, and after three failed attempts in a row the campaign stops itself \gl{D5-3}; the cause is found before it starts again.
  \item \textbf{Kafka runs keep pausing}\chgold{D6-10} for fractions of a second. The pauses are reported with the TCP counters and the broker's log of each run; the analysis uses medians and per-message placement, which a few paused messages barely move.
  \item \textbf{A copy does not match its fingerprints.}\chgold{D5-4} It is made again from the driver. A file that still does not match is reported, and its run is listed as lost, never silently dropped.
\end{itemize}

\section[How this plan is frozen]{How this plan is frozen\protect\chgold{D4-13}}\label{sec:freeze}
A pre-registration is only worth something if a reader can check that it was written before the data. This plan is frozen by committing it in full, its source and this PDF with their SHA-256 fingerprints, to the \texttt{freezes} folder of the project's public repository, before the first Stage 0 run. The code that applies Section~\ref{sec:operations} is in the same repository, and every run records the version of it that ran.\chgold{D5-6} A later change is a new, numbered freeze whose first section lists what changed and why, never an edit of this one. This version is the third freeze \gl{freeze 03}.\chgold{D6-12}\chgold{D7-4} Results are reported against the freeze in force when they ran, including every prediction that failed.

\end{document}
